\documentclass[11pt]{article}
\usepackage[a4paper,margin=25mm]{geometry}
\usepackage{amsmath,amssymb,amsthm,mathtools}
\usepackage{mathrsfs}
\usepackage{booktabs,longtable,array,tabularx}
\usepackage{microtype}
\usepackage{xcolor}
\usepackage[colorlinks=true,linkcolor=blue!60!black,citecolor=blue!60!black,urlcolor=blue!65!black]{hyperref}
\usepackage[nameinlink,noabbrev]{cleveref}
\setlength{\emergencystretch}{3em}
\newtheorem{theorem}{Theorem}[section]
\newtheorem{lemma}[theorem]{Lemma}
\newtheorem{proposition}[theorem]{Proposition}
\newtheorem{corollary}[theorem]{Corollary}
\theoremstyle{definition}
\newtheorem{definition}[theorem]{Definition}
\newtheorem{remark}[theorem]{Remark}
\newcommand{\A}{\mathcal A}
\newcommand{\T}{\mathcal T}
\newcommand{\B}{\mathcal B}
\newcommand{\ind}{\mathbf 1}
\newcommand{\eps}{\varepsilon}
\newcommand{\Th}{\Theta(N)}
\newcommand{\phistar}{\varphi^{*}}
\newcommand{\Fup}{\mathcal F}
\newcommand{\Flow}{\mathcal L}
\newcommand{\cert}[1]{\texttt{\detokenize{#1}}}
\newcommand{\statusclosed}{\textsf{CLOSED}}
\title{Positive Chen representations for $\log\log N\ge24.9$\\
\large A finite Wu--BJS proof and its reproducible audit}
\author{Cheng Huang\\
\small Shanghai Jiaotong University}
\date{5 September 2026}

\begin{document}
\maketitle
\begin{abstract}
The objective is positivity, not preservation of a proposed numerical
coefficient at $23.36$.  We use the published product cutoff $10^9$, a
fixed-cutoff prime--integer rectangle estimate proved below in the required
range $X\ge N^{0.398}$, and a complementary-prime sieve of level $N^{2/5}$.
The latter retains the original $p_2$-roughness, which makes both grouped
coefficients bounded by one.  Its small-conductor contribution is estimated
before Perron separation; its large-conductor contribution is estimated by
the multiplicative large sieve.  All exceptional alternatives are covered
by common worst-case envelopes.  The arithmetic certificate uses Arb balls;
the extension from one endpoint to a half-line is proved by analytic
majorants, not inferred from endpoint sampling.
\end{abstract}
\tableofcontents

\section{Statement, notation and audit scope}

For even $N$, let
\[
 D_{1,2}(N)=\#\{p<N:p\text{ is prime},\ \Omega(N-p)\le2\}.
\]
We count prime factors with multiplicity.  Put
\[
 L=\log N,\quad c=\log L,\quad N_0=\exp(\exp(249/10)),
\]
\[
 C_N=\prod_{p>2}\left(1-\frac1{(p-1)^2}\right)
       \prod_{\substack{p\mid N\\p>2}}\frac{p-1}{p-2},
 \qquad \Th=\frac{C_NN}{L^2}.
\]
The elementary bound $C_N>1/2$ suffices throughout: deleting factors from
the convergent product
\[
 \prod_{n=2}^{\infty}(1-n^{-2})=\frac12
\]
increases it, and the factors indexed by $n=p-1$ form a proper subset.
We also use the elementary upper bound $C_N<1+\omega(N)$ only for
finite exclusions: if $q_1<\cdots<q_j$ are the odd primes dividing
$N$, then $q_i\ge i+2$ and
$\prod_i(q_i-1)/(q_i-2)\le\prod_{i=1}^j(i+1)/i=j+1$.

\begin{theorem}\label{thm:main}
Every even integer $N\ge N_0$ satisfies $D_{1,2}(N)>0$.
\end{theorem}

The proof is given in \cref{sec:assembly,sec:uniform}.  No assertion at
$c=23.36$ is retained.  The supplied \cert{reply.md} is used as the audit
checklist, not as a mathematical source.  References to BJS mean exactly
arXiv:2207.09452v6 \cite{BJS}; references to Wu mean Lemmas 9.1--9.2 of
\cite{Wu}.  New statements needed outside their printed ranges are proved
here and are not attributed to those papers.

All prime intervals are left-closed and right-open unless an inequality
explicitly says otherwise.  A finite sequence is a multiset: equal integer
values with different source tags retain their multiplicities.  Set
\[
 P_M(t)=\prod_{\substack{p<t\\p\nmid M}}p,
 \qquad P_{\rm all}(t)=\prod_{p<t}p,
\]
\[
 \A=\{N-p:p<N\text{ prime},\ p\nmid N\},\qquad
 \A_a=\{b/a:b\in\A,\ a\mid b\},
\]
and write $S(\A_a,P_M(t))$ for the number of members of $\A_a$ coprime to
$P_M(t)$.  Thus every member of $\A$ is odd and coprime to $N$.

The retuned rational parameters are
\[
\begin{aligned}
 \kappa_1&=\frac{191}{2000},&
 \kappa_2&=\frac{383}{4000},&
 \rho_0&=\frac{119}{500},&
 \sigma_2&=\frac{6047}{20000},\\
 \sigma_1&=\frac{121}{400},&
 \delta&=\frac{29}{25},&
 \lambda&=\frac12-2\kappa_1=\frac{309}{1000}.&&
\end{aligned}
\]
\[
 z_*=N^{\kappa_1},\quad y=N^{1/3},\quad z=N^{1/8},\quad
 u_0=10^9,\quad \eps_0=\frac1{10000},\quad
 \eta=\frac{1452}{10^{10}}.
\]
The parameters called $\eps,Z,z_0,H$ and the Johnston terminal in the
earlier manuscript are not used.  In particular, there is no undefined
support parameter: the complementary level is exactly $N^{2/5}$.

\section{External inputs, products and model masses}\label{sec:inputs}

\subsection{The precise linear-sieve input}

Let $F,f$ be the functions in BJS Theorem 6:
\[
 F(s)=2e^\gamma/s,\quad f(s)=0\quad(0<s\le2),\qquad
 (sF(s))'=f(s-1),\quad(sf(s))'=F(s-1)\quad(s\ge2).
\]
On $s\ge1$, $F\ge1$ is decreasing, and $0\le f\le1$ is increasing.
These properties also follow from the nonnegative decreasing functions
$f_n$ in BJS equations defining $F=1+\sum_{n\ {\rm odd}}f_n$ and
$f=1-\sum_{n\ {\rm even}}f_n$.  Define
\[
 h(s)=\begin{cases}e^{-2},&1\le s\le2,\\e^{-s},&2\le s\le3,\\
 3e^{-s}/s,&s\ge3,\end{cases}\qquad
 \Fup(s)=F(s)+106\eta e^2h(s),\quad
 f_-(s)=f(s)-107\eta e^2h(s).
\]
The constants $106,107$ are the outward integer bounds in BJS Table 1
at reciprocal tolerance $10^5$.  They remain valid for our smaller
tolerance, as stated and proved there.

BJS Theorem 6 uses a tag-dependent multiplicative function $g_\alpha$,
with $0\le g_\alpha(p)<1$.  Its main term is
\[
 X_A=\sum_\alpha w_\alpha\prod_{p\mid P(w)}(1-g_\alpha(p)),
\]
not an aggregate multiplicative density.  The product hypothesis is
\begin{equation}\label{eq:product-hyp}
 \prod_{\substack{p\in\mathbb P\setminus\mathbb Q\\u\le p<w}}
 (1-g_\alpha(p))^{-1}<(1+\eta)\frac{\log w}{\log u}
 \quad\text{for every tag and every }1<u<w.
\end{equation}
Here $Q=\prod_{p\in\mathbb Q}p$.  The upper inequality requires $D\ge w$,
the lower inequality requires $D\ge w^2$, and the remainder is the sum
of absolute discrepancies over $d\mid P(w),d<QD$.

We also use the following model-mass formulation.  It records explicitly
the minor extension from the printed formulation with actual total mass.

\begin{lemma}[Finite tagged model sieve]\label{lem:model}
Suppose finite nonnegative model masses $x_\alpha$ and multiplicative
$g_\alpha$ satisfy \eqref{eq:product-hyp}.  Suppose the actual divisibility
counts have the form
\[
 A_d=\sum_\alpha x_\alpha g_\alpha(d)+r(d).
\]
Then, with $V_\alpha(w)=\prod_{p\mid P(w)}(1-g_\alpha(p))$ and
$s=\log D/\log w$,
\[
 S(A,P(w))\le\Fup(s)\sum_\alpha x_\alpha V_\alpha(w)
               +\sum_{\substack{d\mid P(w)\\d<QD}}|r(d)|
 \quad(D\ge w),
\]
\[
 S(A,P(w))\ge f_-(s)\sum_\alpha x_\alpha V_\alpha(w)
               -\sum_{\substack{d\mid P(w)\\d<QD}}|r(d)|
 \quad(D\ge w^2).
\]
The discrepancy at $d=1$ is included if the model total is not the actual
total.  There is no requirement that the weighted average of the
$g_\alpha$ be multiplicative.
\end{lemma}
\begin{proof}
Use the Rosser upper and lower coefficients in the proof of BJS
Theorem 6 (equivalently its Proposition 13), including its exact small-prime
presieve.  They depend on $P(w),D,Q$, not on the model masses or tags;
their absolute values are at most one and their support is $d<QD$.
Their divisor sums bound $\ind_{(n,P(w))=1}$ pointwise.  Finite
distributivity gives
\[
 \sum_d\lambda^\pm(d)A_d
 =\sum_\alpha x_\alpha\sum_d\lambda^\pm(d)g_\alpha(d)
  +\sum_d\lambda^\pm(d)r(d).
\]
BJS bounds each inner model sum by $\Fup(s)V_\alpha(w)$ or
$f_-(s)V_\alpha(w)$ under \eqref{eq:product-hyp}.  Taking the absolute
value only of the final discrepancy sum proves both assertions.  This
also proves the multiset version without identifying unequal tag densities.
\end{proof}

\subsection{The published \texorpdfstring{$10^9$}{1,000,000,000} product cutoff}

\begin{lemma}\label{lem:product}
For every sieve used here, choose $\mathbb Q=\{p\in\mathbb P:p\le10^9\}$.
Then \eqref{eq:product-hyp} holds with the stated $\eta$, and
\begin{equation}\label{eq:Qbound}
 \log Q<\overline q:=1000000050.
\end{equation}
This includes densities $1/(p-1)$, densities obtained by setting any of
them to zero, and the integer density $1/p$.
\end{lemma}
\begin{proof}
BJS Lemma 18 states, for $w>\exp4000$ and $10^9<u<w$,
\[
 \prod_{u\le p<w}(1-1/(p-1))^{-1}
 <(1+\eta)\log w/\log u.
\]
All our sift limits satisfy $\log w\ge\kappa_1e^{24.9}>4000$.
Deleting primes or decreasing their densities decreases this product.
If $1<u\le10^9$, choose $v>10^9$ below both $w$ and the least prime
above $10^9$.  The primes outside $\mathbb Q$ are unchanged by replacing
the lower endpoint $u$ by $v$.  The published inequality at $v$ has a
smaller right side than that required at $u$.  This proves the full
quantifier in \eqref{eq:product-hyp}; the empty product case is immediate.
Finally BJS Lemma 25 gives
\[
 \log Q\le\vartheta(10^9)
 <10^9\left(1+\frac{9\cdot10^{-7}}{\log10^9}\right)
 <1000000050.
\]
The last strict inequality is checked with Arb in the structural
certificate.  Thus the unproved $10^7$ product lemma is replaced, not
silently assumed.
\end{proof}

Set
\begin{equation}\label{eq:levels}
 \theta=\frac{2423257}{5000000},\qquad
 \Delta=\frac{81}{10^9},\qquad \theta_- =\theta-\kappa_1\Delta,
 \qquad H_N=\frac{\sqrt{x_1(N)}}{\log^{10}x_1(N)},
 \quad x_1(N)=N/L^5.
\end{equation}
The certificate verifies at $c=24.9$ that
\[
 \theta<\frac12-\frac{\overline q+20c}{e^c},\qquad
 \frac{20c}{\kappa_1e^c}<\Delta.
\]
The first right side increases and the second decreases on $c\ge24.9$.
Consequently, on the entire half-line,
\begin{equation}\label{eq:level-support}
 QN^\theta\le \sqrt N/L^{20}<H_N,
 \qquad \frac{20c}{\log w}\le\Delta\quad(w\ge z_*).
\end{equation}
Indeed $\log H_N\ge L/2-12.5c>L/2-20c$.

\subsection{Prime-counting and arithmetic-progression inputs}

Here are the remaining external inputs, with the exact conditions needed.
BJS Lemmas 17, 24, 26, 29, 31, 34, 35 and 42--43 are used; the names
of their formulae below remove any ambiguity about numbering.
\begin{enumerate}
\item For $x>\exp4000$, the reciprocal-prime sum differs from
$\log\log x+M$ by a quantity in
$[-2.964\cdot10^{-6}/\log x,\ 1.436\cdot10^{-16}/\log x]$.
Also $\sum_{d\le t}\mu^2(d)/\varphi(d)\le1.1\log t$ for $t\ge10^9$.
\item The explicit prime-counting inequalities used in BJS give
\begin{equation}\label{eq:A-mass}
 N/L<|\A|<1.00005N/L,
 \qquad \pi(x)<x/(\log x-3/2)
\end{equation}
in all ranges used below (every such $\log x$ is larger than $4000$).
The hypotheses $\log\log x_1(N)\ge10.4$ are much weaker than $c\ge24.9$.
\item For $x_2(T)=T/\log^{15}T$, $T>X$, and
$r<\log^{10}x_2(T)$, BJS Lemma 29 gives
\begin{equation}\label{eq:chi-prime}
 \sum_{\substack{\chi\bmod r\\\chi\ne\chi_0}}|\pi(T,\chi)|
 <v_r(X)T/\log^5T,
 \quad \log\log x_2(X)\ge10.4.
\end{equation}
Its displayed formula is reproduced in \cref{eq:vformula} below.
\item Let $k$ be the exceptional conductor up to
$Q_1(x_1(N))=\log^{10}x_1(N)$, setting $k=0$ if it is absent or not
coprime to $N$.  A nonzero relevant conductor is odd and square-free.
If $k<\log^\delta x_1(N)$, BJS's small-exceptional AP estimate applies
to all square-free $d\le H_N$ coprime to $N$.  In the other case its
nonexceptional estimate applies to those $d$ with $k\nmid d$, and its
relative estimate applies, with this fixed full conductor $k$, to
\[
 \pi(N;kd,N)-\pi(N;k,N)/\varphi(d),\quad
 d\le H_N/k,\quad(d,kN)=1.
\]
We use the relative estimate only for $k$ itself, not for an arbitrary
additional large source tag.
\end{enumerate}

In the notation $r(d)=|\A_d|-|\A|/\varphi(d)$ and
$r_k(d)=|\A_{kd}|-|\A_k|/\varphi(d)$, these AP estimates imply
\begin{equation}\label{eq:AP}
 \sum |r(d)|\le b(N)N/L^3,\qquad
 \sum |r_k(d)|\le b(N)N/L^3,\qquad b(N)<2.21/L^2,
\end{equation}
over their respective ranges.  For the first sum in the large case,
$k\nmid d$ is imposed.  The coefficient $b(N)$ is the BJS $c_4$ formula
with a weaker zero bound, reproduced in \cref{sec:APformula}.  Passing
from $\pi$ to $\A$ costs at most $\omega(N)$ per modulus; the
$c_4$ term already includes this.  Neither this passage nor the relative
estimate requires that all prime factors of $d$ be below a particular
sift limit.  This is also immediate from their proofs, which sum over
all square-free moduli.  All scale and zero-bound conditions are verified
in \cref{sec:uniform} and in the structural certificate.

For the remaining elementary external inputs we use exactly the forms
quoted and used by BJS:
\begin{itemize}
\item Brun--Titchmarsh:
$\pi(x;q,a)<2x/(\varphi(q)\log(x/q))$ for $(a,q)=1$, $x>q$.
For its direct uses here the moduli are at most $L^{10}$ or $H_N$,
and the arguments are at least $x_1(N)$, so $x>q$ holds.
The smaller endpoints internal to BJS Lemma 29 are covered by that
lemma, whose separate scale hypothesis is verified below.
\item Rosser--Schoenfeld:
$q/\varphi(q)<e^\gamma\log\log q+2.5/\log\log q$ for $q\ge3$.
We use it only for an odd exceptional conductor whose
$\log\log q>2$.
\item The full Mertens product estimate used in BJS's proof of
Lemma 38 has relative error at most $1/(2\log^2x)$ for $x\ge285$;
using $1/\log^2x$ permits either endpoint convention.
\item The multiplicative large sieve in the proof of BJS Lemma 32:
for any complex coefficients on a prefix of length $M$,
\[
 \sum_{r\le R}\frac r{\varphi(r)}
       \sum_{\chi\bmod r}^{*}\left|\sum_{n\le M}a_n\chi(n)\right|^2
 \le (M+3R^2)\sum_{n\le M}|a_n|^2 .
\]
Here the star restricts to primitive characters.  This inequality
requires neither $R^2\le M$ nor a prime condition on the coefficients.
For $r\in[R,2R)$ it yields the multiplier $M+12R^2$ used below.
\end{itemize}
These four inputs require no distribution assertion beyond their
displayed ranges.  The new rectangular and switched applications are
proved below, rather than quoted as further external theorems.

For reference, the complete external-input condition ledger is below.
It is a logical ledger: a row names every hypothesis used later and the
place where it is discharged.
\begin{center}
\small
\begin{tabularx}{\textwidth}{@{}>{\raggedright\arraybackslash}p{0.19\textwidth}
  >{\raggedright\arraybackslash}X>{\raggedright\arraybackslash}X@{}}
\toprule
Input & Hypotheses actually used & Discharge in this proof\\
\midrule
Wu Lemmas 9.1--9.2
 & $0<\kappa_1<\kappa_2<\rho_0<\sigma_2<\sigma_1<1/3$ and
   $3\kappa_1+\rho_0\ge1/2$
 & Exact rational checks precede \cref{lem:wu}; the finite pointwise and
   Buchstab proof is repeated there, so Wu's asymptotic $O$-term is not
   imported.\\
BJS Theorem 6
 & Nonnegative finite sequence, $0\le g_\alpha(p)<1$, the product
   condition for every tag and $1<u<w$, $D\ge w$ (upper) or $D\ge w^2$
   (lower), and discrepancies for $d<QD$
 & \Cref{lem:model,lem:product}; every later application states its
   level and remainder support.  The tag masses are not replaced by an
   aggregate density.\\
BJS product and Chebyshev bounds
 & $w>\exp4000$, $10^9<u<w$, and the product $Q$ of the presieving primes
 & \Cref{lem:product,eq:level-support}; the cases $u\le10^9$ and an empty
   product are proved separately.\\
BJS Lemma 29
 & Endpoint $T>X$, $\log\log x_2(X)\ge10.4$, and
   $r<\log^{10}x_2(T)$, with the applicable exceptional-zero bound
 & \Cref{eq:vformula,thm:rectangle,thm:switch}; both endpoints are kept,
   and the fixed cutoffs are strictly below the required powers.\\
BJS small-modulus AP estimate
 & $k<\log^\delta x_1(N)$, square-free $d\le H_N$, $(d,N)=1$
 & \Cref{eq:AP,sec:APformula}; the endpoint maximum and its half-line
   propagation are proved in \cref{sec:uniform}.\\
BJS nonexceptional AP estimate
 & $k\ge\log^\delta x_1(N)$ and $k\nmid d$
 & The first terms of \eqref{eq:chain-identity}; the missing
   $q_{j+1}$ makes $k\nmid am_jd$ exactly.\\
BJS relative AP estimate
 & One fixed nontrivial full conductor $k$, $(d,kN)=1$, and
   $d\le H_N/k$
 & Only the last term of \eqref{eq:chain-identity}; here $k$ is the full
   exceptional conductor and $ad\le H_N/k$.\\
Multiplicative large sieve
 & Primitive characters and arbitrary finite complex coefficient
   sequences
 & \Cref{thm:rectangle,thm:switch}; the four positive terms after
   Cauchy--Schwarz and every dyadic interval are explicitly summed.\\
\bottomrule
\end{tabularx}
\end{center}

\section{A finite signed Wu reduction}\label{sec:wu}

Define $U_i$ as follows.  All summation variables are primes and all
displayed source products are required to be coprime to $N$.
For brevity $S_M(a,w)=S(\A_a,P_M(w))$.
\begin{align*}
 U_1&=\sum_{z_*\le p<N^{\kappa_2}}S_N(p,p),&
 U_2&=\sum_{z_*\le p<N^{\sigma_1}}S_N(p,z_*),&
 U_3&=\sum_{z_*\le p<N^{\sigma_2}}S_N(p,z_*),\\
 U_4&=\sum_{z_*\le p_1<p_2<N^{\kappa_2}}S_N(p_1p_2,z_*),\\
 U_5&=\sum_{z_*\le p_1<N^{\kappa_2}\le p_2<N^{\lambda}/p_1}
 S_N(p_1p_2,z_*),\\
 U_6&=\sum_{z_*\le p_1<N^{\kappa_2},\,N^{\rho_0}\le p_2<N^{\sigma_2}}
 S_N(p_1p_2,p_1).
\end{align*}
For $i=7,8$, let
\[
 U_i=\sum_{N^{\tau_i}\le p_1<p_2<\sqrt{N/p_1}}S_{Np_1}(p_1p_2,p_2),
 \quad\tau_7=\sigma_1,\quad\tau_8=\sigma_2.
\]
The remaining rows are
\begin{align*}
 U_9&=\sum_{z_*\le p_1<N^{\sigma_1}\le p_2<\sqrt{N/p_1}}
 S_{Np_1}(p_1p_2,p_2),\\
 U_{10}&=\sum_{N^{\kappa_2}\le p_1<N^{\sigma_2}\le p_2<\sqrt{N/p_1}}
 S_{Np_1}(p_1p_2,N^{\sigma_1}),\\
 U_{11}&=\sum_{N^{\kappa_2}\le p_1<p_2<p_3<N^{\sigma_2}}
 S_N(p_1p_2p_3,p_2),\\
 U_{12}&=\sum_{\substack{N^{\kappa_2}\le p_1<N^{\sigma_2}\\
 N^{\sigma_1}\le p_2<p_3<\sqrt{N/p_1}}}S_{Np_1}(p_1p_2p_3,p_2),\\
 U_{13}&=\sum_{z_*\le p_1<p_2<p_3<p_4<N^{\kappa_2}}
 S_N(p_1p_2p_3p_4,p_2),\\
 U_{14}&=\sum_{z_*\le p_1<p_2<p_3<N^{\kappa_2}\le p_4<N^{\lambda}/p_3}
 S_N(p_1p_2p_3p_4,p_2),\\
 U_{15}&=\sum_{\substack{N^{\kappa_2}\le p_1<N^{\sigma_2}\\
 N^{\sigma_2}\le p_2<p_3<p_4<N^{\sigma_1}}}S_N(p_1p_2p_3p_4,p_3).
\end{align*}

\begin{lemma}\label{lem:wu}
For $N\ge N_0$,
\begin{equation}\label{eq:Wu-finite}
\begin{split}
4D_{1,2}(N)\ge{}&4S(\A,P_N(z_*))-U_1-U_2-U_3+U_4+U_5\\
 &-2U_7-2U_8-U_9-U_{10}-U_{13}-U_{14}-E_W,
\end{split}
\end{equation}
where, with $W=L/\log2$ and $B=2N^{11/12}+\sqrt N+W+1$,
$E_W\le(4+W^2)B$.
\end{lemma}
\begin{proof}
All of Wu's parameter inequalities hold; in particular
\[
 3\kappa_1+\rho_0=\frac{1049}{2000}>\frac12.
\]
This is Wu's Lemma 9.2 condition and is also the exact geometric
inequality used below to truncate the positive double sum.
First restrict to square-free source members coprime to $N$ and sifted
at $z_*$.  Wu's Lemma 9.1 is a pointwise inequality on these members.
Here is a direct verification, using our quotient definition of $\A_d$.
For any $0<\kappa<\sigma<1/3$ and a square-free, $N^\kappa$-rough
source member $a$ with at least three prime factors, let $j$ be its
number of prime factors below $N^\sigma$.
Define $s_1(a)$ to count its prime divisors in $[N^\kappa,N^\sigma)$.
Define $s_2(a)$ to count pairs $p_1<p_2$ dividing $a$ with
$N^\sigma\le p_1$, $p_2<\sqrt{N/p_1}$, and all prime factors of
$a/(p_1p_2)$ at least $p_2$.  Define $s_3(a)$ by the same
conditions except $N^\kappa\le p_1<N^\sigma\le p_2$.
Finally $s_4(a)$ counts triples
$N^\kappa\le p_1<p_2<p_3<N^\sigma$ dividing $a$ for which all
prime factors of $a/(p_1p_2p_3)$ are at least $p_2$.
All selected factors are coprime to $N$, since $a$ is.
The four counts in the pointwise weight
$1-s_1(a)/2-s_2(a)-s_3(a)/2+s_4(a)/2$ are therefore
\[
\begin{array}{c|cccc}
 j&s_1&s_2&s_3&s_4\\ \hline
 0&0&1&0&0\\
 1&1&0&1&0\\
 2&2&0&0&0\\
 j\ge3&j&0&0&j-2
\end{array}
\]
Indeed a sift at the second selected prime forces the first two
selected factors to be the least two factors of $a$; the optional
third selected factor has $j-2$ choices.  The restriction
$p_2<\sqrt{N/p_1}$ is automatic for these least two factors, since
$a\ge p_1p_2p_3>p_1p_2^2$ and $a<N$.
Every row of the table gives weight zero.  If $\Omega(a)\le2$,
$s_4(a)=0$ and the weight is at most one.  This proves the needed
pointwise inequality with no asymptotic error.

The finite Buchstab identity used throughout is
\[
 S(\A_d,P_M(v))=S(\A_d,P_M(u))
 -\sum_{\substack{u\le q<v\\q\nmid M}}S(\A_{dq},P_M(q)).
\]
It assigns a residual failing the larger sift to its least prime in
$[u,v)$.  Repeated selected primes contribute zero on the square-free
source.  Deleting a selected prime from $P_M$ has no effect on such a
residual.  These observations explain all $P_N$ versus $P_{Np_1}$
uses in the following finite algebra.

For clarity, the finite algebra in Lemma 9.2 has the following three
stages.  Applying Lemma 9.1 with $(\kappa_2,\sigma_2)$ and using
Buchstab's identity gives
\[
 2D\ge 2S-U_1-U_3+U_4+U_5+U_6-2U_8-U_{10}
             +U_{11}+U_{12}+\Delta_1.
\]
Here $\Delta_1$ is the negative sum of the three triple-prime domains
\[
\begin{array}{ll}
 z_*\le p_1<p_2<p_3<N^{\kappa_2},&\text{sift at }p_1,\\
 z_*\le p_1<p_2<N^{\kappa_2}\le p_3<N^{\lambda}/p_2,
 &\text{sift at }p_1,\\
 N^{\kappa_2}\le p_1<N^{\sigma_2}\le p_2<p_3<N^{\sigma_1},
 &\text{sift at }p_3.
\end{array}
\]
The first two use $P_N$ and the last uses $P_{Np_1}$.  The truncation
of the positive double sum at $N^{\lambda}/p_1$ is valid because
$p_1\ge N^{\kappa_1}$ and $3\kappa_1+\rho_0>1/2$, whence
$N^{\rho_0}\ge N^{\lambda}/p_1$.  The second use of Lemma 9.1, with
$(\kappa_1,\sigma_1)$, gives
\[
 2D\ge2S-U_2-2U_7-U_9+T,
 \quad T=\sum_{z_*\le p_1<p_2<p_3<N^{\sigma_1}}S_N(p_1p_2p_3,p_2).
\]
Each of the three domains above is contained in the domain of $T$.
Subtracting their sift limits by the finite Buchstab identity gives
$\Delta_1+T\ge-U_{13}-U_{14}+U_{15}$.
Explicitly, in the first two domains the additional prime $q$ obeys
$p_1<q<p_2$: relabel $(p_1,q,p_2,p_3)$ as the four increasing
primes.  The new residual is sifted at $q$, the second prime, and
in the second domain the old bound $p_3<N^{\lambda}/p_2$ becomes
exactly the bound defining $U_{14}$.  In the third domain,
$p_2<q<p_3$ and the relabeling $(p_1,p_2,q,p_3)$ gives $U_{15}$
with sift at its third prime.  The endpoints $q=p_1$ or $q=p_2$
vanish by square-freeness; the parts of $T$ outside the three
disjoint domains are nonnegative.
This proves the fifteen-row inequality on the restricted source.  All
square-divisor error terms in Wu's printed asymptotic statement vanish
on this square-free source.

Delete the four nonnegative rows $U_6,U_{11},U_{12},U_{15}$ at this
point.  A sifted source failing square-freeness has $q^2\mid a$ for
some $q\ge z_*$, and their number is bounded by
\[
 \sum_{z_*\le q\le\sqrt N}(N/q^2+1)
 \le2N/z_*+\sqrt N\le B.
\]
The congruence floor term is retained.  Each bad member contributes at
most $4+2\binom{\omega(a)}2\le4+W^2$ to the retained positive rows.
Extending negative rows to all of $\A$ only lowers the right side.
This proves \eqref{eq:Wu-finite}.  In particular no three- or four-prime
positive multiplicity is charged with a two-prime bound.
\end{proof}

\section{Ordinary sieve rows and the exceptional chain}\label{sec:ordinary}

\subsection{A common lower envelope}

Set $e_1=1/35$ and, for $u\ge1$, define
\begin{equation}\label{eq:lower-envelope}
 \Flow(u)=\begin{cases}
 -4e_1-(1+4e_1)107\eta e^2e^{-2},&1\le u<2,\\
 \begin{aligned}
 &f(u)-107\eta e^2h(u)\\
 &\quad-3e_1\{F(u)-f(u)+213\eta e^2h(u)\},
 \end{aligned}&u\ge2.
 \end{cases}
\end{equation}
The first branch is negative and is used only as a trivial lower bound.
On the second branch $f_-$ is increasing and $\Fup$ is decreasing;
hence $f_--3e_1(\Fup-f_-)$ is increasing and does not exceed $f_-$.
This uses $F-f$ itself, rather than an uncertified numerical maximum.

\begin{lemma}[Exceptional chain with an exact final mass]\label{lem:exceptional}
The lower ordinary rows may use \eqref{eq:lower-envelope} at
$u=(\theta-\log a/L)/\kappa_1-\Delta$ in both exceptional alternatives.
All upper rows may use a factor $36/35$ and their indicated, possibly
smaller, sieve arguments.  The normalized finite cost of the exceptional
main-mass replacement is at most
\begin{equation}\label{eq:exc-cost}
 100 B_E(N)+1000/L^2<10^{-6},
 \quad B_E(N)=\frac{96\log^2(10c)}{\log^\delta x_1(N)}.
\end{equation}
AP discrepancies are charged separately in \eqref{eq:ordinary-AP-cost}.
\end{lemma}
\begin{proof}
In the small alternative use \cref{lem:model} and the all-modulus
estimate \eqref{eq:AP}.  For a nonnegative lower model expression,
$|\A|>N/L$ may be used; for a negative expression the same lower
bound follows instead from the nonnegativity of the actual sifted count.
Since $\Flow\le f_-$, this gives the stated lower envelope.

In the large alternative write the relevant full conductor as
$k=q_1\cdots q_\ell$ with $q_1>\cdots>q_\ell$ odd primes, and put
$m_j=q_1\cdots q_j$, $m_0=1$, $P^{(j)}(w)=P_N(w)/\prod_{i\le j}q_i$.
Every $q_i<w$, because $k\le L^{10}<z_*$.  The source tags $a$ here
have at most two prime factors, all at least $z_*$, so $(a,k)=1$.
Finite inclusion--exclusion gives, exactly,
\begin{equation}\label{eq:chain-identity}
 S(\A_a,P_N(w))=
 \sum_{j=0}^{\ell-1}(-1)^jS(\A_{am_j},P^{(j+1)}(w))
 +(-1)^\ell S(\A_{ak},P^{(\ell)}(w)).
\end{equation}
This follows by assigning any member not sifted by $P_N(w)$ to the
first prime in the ordered list $q_1,\ldots,q_\ell$ which fails its
coprimality test; the alternating prefix terms cancel.

Use the common level $D_a=N^\theta/(ak^2)$ in every term.  The first
$\ell$ model masses are $|\A|/\varphi(am_j)$; their discrepancies
are $r(am_jd)$.  The last model mass is
$|\A_k|/\varphi(a)$; its discrepancies are $r_k(ad)$.  No approximation
of the full-conductor prime count is used in applying the relative
distribution theorem.  In the first terms $k\nmid am_jd$ because
$q_{j+1}\nmid d$.  In the last term $ad\le H_N/k$ by
\eqref{eq:level-support}.  All combined moduli are square-free and
coprime to $N$.  These are precisely the two AP ranges in \eqref{eq:AP}.
Moreover
\[
 \frac{\log D_a}{\log w}\ge
 \frac{\theta L-\log a}{\log w}-\Delta.
\]
If the right side is below $2$, its negative lower envelope follows
from $S\ge0$; no lower linear sieve is applied below its allowed level.

Let $V_j=\prod_{p\mid P^{(j)}(w)}(1-1/(p-1))$ and
$\phistar(m_j)=\prod_{i\le j}(q_i-2)$.  The ideal model coefficients
satisfy
\[
 \frac{V_{j+1}}{\varphi(m_j)V_0}
 =\frac1{\phistar(m_j)}\left(1+\frac1{q_{j+1}-2}\right),\qquad
 \frac{V_\ell}{\varphi(k)V_0}=\frac1{\phistar(k)}.
\]
Thus their alternating sum in \eqref{eq:chain-identity} is exactly one.
Their sum for $j\ge1$ is at most
\[
 \frac1{q_1-2}\left(1+\frac2{q_2-2}
 +\frac2{(q_2-2)(q_3-2)}+\cdots\right)\le\frac3{q_1-2}.
\]
For completeness the tail without the initial one is at most $2$:
each of its $\ell-1$ denominator products is at least $q_2-2$,
and $q_2\ge2\ell-1$, giving
$(\ell-1)/(2\ell-3)\le1$.  The case $\ell=1$ has no tail.
The fixed primorial comparison
\[
 \prod_{2<p\le31}p=100280245065
 <\log^\delta x_1(N_0)
\]
implies $q_1\ge37$ for every subsequent $N$: an odd square-free
integer all of whose prime factors are at most $31$ is at most the
displayed primorial.  The $j=0$ local factor is therefore at most
$1+e_1=36/35$.  Let $P_+$ and $P_-$ be the sums of the ideal
coefficients of the positive and negative terms in
\eqref{eq:chain-identity}.  The exact identity above gives
$P_+-P_-=1$, while the tail estimate gives $P_-\le3e_1$.  Hence
\[
 f_-P_+-\Fup P_-
 =f_- -P_-(\Fup-f_-)
 \ge f_- -3e_1(\Fup-f_-)=\Flow.
\]
This accounts explicitly for both the $106\eta$ upper-sieve correction
and the $107\eta$ lower-sieve correction.  Applying lower weights to
the positive terms and upper weights to the negative terms now gives
\eqref{eq:lower-envelope}, apart
from the final-mass defect and the separately charged AP errors.

The Brun--Titchmarsh inequality and \eqref{eq:A-mass} give
\[
 0\le\frac{|\A_k|}{|\A|}\le
 \frac{2L}{\varphi(k)(L-\log k)}<\frac3{\varphi(k)}.
\]
Hence the defect, after multiplication by $V_\ell/V_0$, is at most
$3/\phistar(k)$.  The Rosser--Schoenfeld bound for $k/\varphi(k)$
used by BJS, together with $\log^\delta x_1(N)\le k\le L^{10}$,
gives $k/\varphi(k)\le4\log(10c)$ here:
if $t=\log\log k>2$, then
$e^\gamma t+2.5/t<(e^\gamma+5/8)t<4t$ and
$t\le\log(10c)$.  Also
\[
 \frac1{\phistar(k)}
 =\frac1k\left(\frac{k}{\varphi(k)}\right)^2
    \prod_{p\mid k}\frac{(p-1)^2}{p(p-2)}
 <\frac{32\log^2(10c)}{k}.
\]
This proves the bound $B_E$ for the defect and for the additional
tail-density cost.  The reciprocal-prime mass of
$[N^{\kappa_1},N^{\kappa_2})$ is less than $1/3$ by
\eqref{eq:measure-discrepancy}; the total available mass is below $3$.
Each of the two positive pair domains therefore has reciprocal-prime
mass below one; $2e^{-\gamma}/\kappa_1<14$; and
$\Fup(u)<2$ for $u\ge2$.  Including the base row of weight four
and both pair rows, the defect multiplier after division by four is
less than $(4+1+1)\cdot14\cdot2/4=42<100$.  The local-product
endpoints fit the additional $1000/L^2$ in \eqref{eq:exc-cost}.
For an upper row, omit $q_1$ from its prime set.  This only enlarges the
sifted set, multiplies its model product by at most $36/35$, and removes
all multiples of $k$ from its remainder support.  This proves the
upper assertion as well.
\end{proof}

\subsection{Row formulae and the AP multiplicity}

Define the universal lower quantities
\begin{align}
 L_0&=\frac{2e^{-\gamma}}{\kappa_1}
       \Flow(\theta/\kappa_1-\Delta),\label{eq:L0}\\
 L_4&=\frac{2e^{-\gamma}}{\kappa_1}
 \int_{\kappa_1}^{\kappa_2}\frac{dt}{t}
 \int_t^{\kappa_2}\Flow((\theta-t-u)/\kappa_1-\Delta)\frac{du}{u},\\
 L_5&=\frac{2e^{-\gamma}}{\kappa_1}
 \int_{\kappa_1}^{\kappa_2}\frac{dt}{t}
 \int_{\kappa_2}^{\lambda-t}\Flow((\theta-t-u)/\kappa_1-\Delta)\frac{du}{u}.
\end{align}
The upper quantities, before the common exceptional factor, are
\begin{align}
 A_1&=2e^{-\gamma}\int_{\kappa_1}^{\kappa_2}
 \Fup((\theta_--t)/t)\frac{dt}{t^2},\\
 A_2&=\frac{2e^{-\gamma}}{\kappa_1}
 \int_{\kappa_1}^{\sigma_1}\Fup((\theta_--t)/\kappa_1)\frac{dt}{t},\\
 A_3&=\frac{2e^{-\gamma}}{\kappa_1}
 \int_{\kappa_1}^{\sigma_2}\Fup((\theta_--t)/\kappa_1)\frac{dt}{t}.
\end{align}
For $S(\A_a,P_N(w))$ the ordinary level is $N^\theta/a$.  All primes
of $a$ are at least $w$, so $(a,d)=1$ on the remainder support.
The model mass is $|\A|/\varphi(a)$ and its discrepancy is exactly
$r(ad)$ in the small or nonexceptional case.  Thus these formulae
follow from \cref{lem:model,lem:exceptional} after the prime-sum transfer
proved in \cref{sec:transfer}.  The replacement of $1/\varphi(a)$ by
$1/a$, and removal of primes dividing $N$ from lower integrals, are
included in the local and finite charges there.

There is no factor equal to the number of source primes in the AP error.
For a fixed combined square-free modulus, a tag $a$ of at most two primes
can be selected in at most
$1+\omega(d)+\binom{\omega(d)}2\le(1+L/\log2)^2$ ways.  In the
exceptional prefix terms the prefix length $j<\ell$ is uniquely
determined by the first absent $q_{j+1}$, so it adds no factor $\ell$.
The final prefix uses one fixed relative AP sum.  The weights are
$4$ for the base lower row, $1$ for each of the two positive pair rows,
and $1$ for each of the three upper rows.  At most two AP sums occur for
a lower row, hence the total multiplier is
$2(4+1+1)+1+1+1=15<16$.  Consequently
\begin{equation}\label{eq:ordinary-AP-cost}
 \frac{E_{\rm AP}}{4\Th}\le
 \frac{16b(N)(1+L/\log2)^2}{4(1/2)L}<10^{-6}.
\end{equation}
The structural certificate encloses the left side at $c=24.9$ below
$5.623\cdot10^{-10}$; the all-$c$ argument is in \cref{sec:uniform}.

\section{The fixed-cutoff rectangle theorem}\label{sec:rectangle}

This section supplies the extension not covered by BJS Lemma 32's
printed hypothesis $X>N^{2/3}$.  The cutoff below is fixed for all
boxes at a given $N$, not changed to $\log^{10}x_2(Y)$ in the last line.

For $X_3=N^{1/8}/2$ set $w_3=\log x_2(X_3)$,
$D_0=w_3^{10}$, and
\[
 \nu=\frac{100}{L^{\delta/2}(\delta c)^2},\qquad\beta=1-\nu.
\]
Define, for either $X=X_3$ or $X=z_*/2$, $w_X=\log x_2(X)$ and
\begin{equation}\label{eq:vformula}
\begin{split}
 v'(X)&=\frac{3.2\cdot10^{-8}}{w_X^4}
  +w_X^4\left(\frac{e^{-\nu w_X}}{1-\nu}
           +1.02w_X^{10}e^{-w_X/2}+3w_X^{10}e^{-2w_X/3}\right),\\
 v(X)&=v'(X)\left(1+\frac1{\log^{10}X\,w_X^5}
       +\frac1{(1-6/w_X)\log X}\right)+\frac3{\log X}.
\end{split}
\end{equation}
The BJS maximum over $t\ge x_2(X)$ is bounded by this endpoint value:
with $x=\log t$, the functions $x^{-4}$, $x^4e^{-\nu x}$,
$x^{14}e^{-x/2}$ and $x^{14}e^{-2x/3}$ decrease for
$x\ge w_X$, since $\nu w_X>100$, $w_X>28$.
The bound $\nu$ is weaker than both zero bounds used by BJS: all small
exceptional thresholds are at most $L^\delta$, and
$\nu<1/(2\cdot2.0452\cdot10c)$.  It therefore works simultaneously
for every allowed conductor, after the indicated exceptional deletion.

\begin{theorem}[Required rectangle extension]\label{thm:rectangle}
Suppose $N\ge N_0$ is even,
\[
 N^{0.398}\le X\le N^{7/8},\quad N^{1/8}\le Z<Y\le N,
 \quad N\le XY\le(1+\eps_0)N,
\]
and $|a(n)|\le1$.  Put $D_R=\sqrt{XY}/\log^{10}Y$.
Let $k_R$ be the relevant exceptional conductor up to $D_0$.
Use all square-free $d<D_R$ coprime to $N$ in the small alternative
$k_R<\log^\delta x_2(X_3)$, and only $k_R\nmid d$ in the other
alternative.  Over either set of moduli,
\begin{equation}\label{eq:rectangle-discrepancy}
 \sum_d\left|\sum_{n<X}\sum_{\substack{Z\le p<Y\\np\equiv N\pmod d}}a(n)
 -\frac1{\varphi(d)}\sum_{n<X}\sum_{\substack{Z\le p<Y\\(np,d)=1}}a(n)\right|
 \le E_1+E_2+E_3+E_4+E_5,
\end{equation}
where the five terms are added, and
\begin{align*}
 E_1&=39v(X_3)XY/\log^3Y,&
 E_2&=108X\log^{13}Y/\log\log Y,\\
 E_3&=26XY\log^2Y/w_3^{10},&
 E_4&=88XY\log^2Y(X^{-1/2}+Y^{-1/2}),\\
 E_5&=106XY/\log^9Y.
\end{align*}
\end{theorem}
\begin{proof}
The hypotheses imply $D_R>10^9$, $D_0<D_R$, $D_R<Y^4$, and
$Z>X_3$; they also imply $\log\log x_2(X_3)>10.4$.
For the third inequality use
$\log D_R\le L/2+\log(1+\eps_0)/2-10\log\log Y<L/2\le4\log Y$.
The other inequalities follow from
$\log D_R\ge L/2-10c$ and are certified at $c=24.9$ with the monotone
majorants in \cref{sec:uniform}.  These are the precise consequences
needed from BJS's original restrictions; no lower bound $X>N^{2/3}$
is used below.

Character orthogonality is valid because $(d,N)=1$.  A character modulo
the square-free $d$ has a unique primitive inducing character modulo
$r\mid d$; writing $d=rs$ gives $(r,s)=1$ and exactly one induced
character per such primitive character.  After removal of the principal
character, the left side is at most
\begin{equation}\label{eq:rectangle-primitive}
 \sum_{rs<D_R}\frac{\mu^2(rs)}{\varphi(r)\varphi(s)}
 \sum_{\substack{\chi\bmod r\ {\rm primitive}\\\chi\ne\chi_0}}
 \left|\sum_{\substack{n<X\\(n,s)=1}}a(n)\chi(n)\right|
 \left|\sum_{\substack{Z\le p<Y\\p\nmid s}}\chi(p)\right|.
\end{equation}
All sums may be restricted by the exceptional deletion.  Such restrictions
are dropped only when this increases a nonnegative majorant.

For $r<D_0$, both prime endpoints are larger than $X_3$ and
$r<\log^{10}x_2(X_3)$, so \eqref{eq:chi-prime} applies at both.
The exact half-open endpoints are $\lceil Y\rceil-1$ and
$\lceil Z\rceil-1$; the function $t/\log^5t$ increases here.  Thus
\[
 \sum_\chi^*\left|\sum_{\substack{Z\le p<Y\\p\nmid s}}\chi(p)\right|
 \le 2v(X_3)Y/\log^5Y+
       \varphi(r)\,1.3841\log D_R/\log\log D_R.
\]
The last term retains every prime dividing $s$.  Bound the integer sum by
$X$ and use $\sum_{n\le D_R}\mu^2(n)/\varphi(n)\le1.1\log D_R$.
For $s\ge16$, BJS's bound
$\omega(s)\le1.3841\log s/\log\log s$ is at most the displayed
bound at $D_R$, since $\log s/\log\log s$ increases there.
For $s<16$, $\omega(s)\le2$ is also below that endpoint bound.
The result is at most
\[
 \frac{2.42v(X_3)XY\log^2D_R}{\log^5Y}
 +\frac{1.21\cdot1.3841D_0X\log^3D_R}{\log\log D_R}
 \le E_1+E_2.
\]
Here $D_0\le\log^{10}Y$, $D_R<Y^4$, and
$t^3/\log t$ is increasing for $t=\log D_R$ in this range.

For $r\ge D_0$, split into $R\le r<2R$, $R=2^jD_0$.
Cauchy--Schwarz and the multiplicative large sieve in BJS's proof give
\begin{align*}
 &\sum_{R\le r<2R}\frac1{\varphi(r)}\sum_\chi^*|A(\chi)P(\chi)|\\
 &\quad\le\frac1R\{(X+12R^2)(Y+12R^2)XY\}^{1/2}\\
 &\quad\le XY/D_0+\sqrt{12X}\,Y+\sqrt{12Y}\,X+12R\sqrt{XY}.
\end{align*}
The support restrictions $(n,s)=1$ and $p\nmid s$ reduce their square
norms and require no extra term.  The number of intervals is at most
$\log D_R/\log2$ and their left endpoints sum to at most $2D_R$.
Summing also over $s$ yields
\[
 1.1\log D_R\left\{\frac{\log D_R}{\log2}
 \left(\frac{XY}{w_3^{10}}+\sqrt{12}XY(X^{-1/2}+Y^{-1/2})\right)
 +24D_R\sqrt{XY}\right\}\le E_3+E_4+E_5.
\]
The constants follow from $1.1\cdot16/\log2<26$,
$1.1\cdot16\sqrt{12}/\log2<88$, and $1.1\cdot24\cdot4<106$.
These inequalities include every conductor interval.  They prove the
claimed fixed-cutoff extension, with no invocation of BJS's out-of-range
printed statement.
\end{proof}

For subsequent aggregation define the slightly larger envelope
\begin{equation}\label{eq:mrect}
\begin{split}
 m_N={}&39v(X_3)+\frac{108\log^{16}X_3}{X_3\log\log X_3}
       +\frac{26L^5}{w_3^{10}}\\
 &+88L^5(N^{-0.199}+N^{-1/16})+\frac{106}{\log^6X_3}.
\end{split}
\end{equation}
Then \eqref{eq:rectangle-discrepancy} is at most $m_NXY/\log^3Y$.
The numerator $L^5$ in the third term is essential because $D_0$ is
fixed.  The exponent $0.199$ in the fourth term uses the actual new
range $X\ge N^{0.398}$, not the old $N^{2/3}$ range.  The function
$\log^{16}Y/(Y\log\log Y)$ decreases for $Y\ge X_3$; this justifies
the second term.  The structural certificate gives
$m_{N_0}<1.437\cdot10^{-8}$.

The condition audit for this proof is as follows.  Each row corresponds
to an actual step above; the original condition $X>N^{2/3}$ has no
unlisted use.
\begin{center}
\small
\begin{tabularx}{\textwidth}{@{}>{\raggedright\arraybackslash}p{0.19\textwidth}
  >{\raggedright\arraybackslash}X>{\raggedright\arraybackslash}X@{}}
\toprule
Step & Required condition & Verification in the new range\\
\midrule
Prime endpoints & $Z-1>X_3$; $\log\log x_2(X_3)>10.4$
 & $Z\ge N^{1/8}=2X_3$; certified scale\\
Low conductor & $r<\log^{10}x_2(X_3)$
 & This is the fixed split $r<D_0$\\
Exceptional zero & Conductor cut at $D_0$
 & Relevant large conductor deleted; common $\nu$ otherwise\\
Induction correction & All primes dividing $s$ counted
 & Explicit $\omega(s)$ term retained in $E_2$\\
Large sieve & Arbitrary complex coefficients of norm at most $X,Y$
 & $|a(n)|\le1$; neither $D_R^2\le X$ nor a Type I/II hypothesis is needed\\
Constants in $E_i$ & $D_R<Y^4$, $D_0\le\log^{10}Y$
 & $Y\ge N^{1/8}$ and $w_3<\log Y$\\
Common envelope & $X^{-1/2}\le N^{-0.199}$, $Y^{-1/2}\le N^{-1/16}$
 & Exactly the new lower bounds on $X,Y$\\
All boxes & Count and weight each box
 & The ceiling sum is performed in \eqref{eq:box-error}\\
\bottomrule
\end{tabularx}
\end{center}

\section{Five finite box families}\label{sec:boxes}

We first explain exactly which triple-prime sources suffice for the
negative rows.  Work on the square-free source from \cref{lem:wu}.
In $U_7,U_8$, a composite residual would imply $p_1p_2^3<N$ with
$p_1,p_2\ge N^{\sigma_2}$, contradicting $4\sigma_2>1$.
In $U_9$ it would imply $N^{\kappa_1+3\sigma_1}<N$, and in
$U_{10}$ it would imply $N^{\kappa_2+\sigma_2+2\sigma_1}<N$.
Both exponent sums exceed one.  Consequently the residual is either
one or a single prime.  A residual equal to one is charged in
\cref{eq:finite-envelope}.  In $U_7,U_8,U_9$ the third prime is at least
$p_2$; in $U_{10}$ it is at least $N^{\sigma_1}$, but need not be
larger than $p_2$.  We do not impose that extra ordering.

For each row in the following table choose the indicated selected prime
$p$ and the complementary two-prime integer $n$.
\begin{center}
\begin{tabular}{@{}ccccc@{}}
\toprule
Family & $p$ & exponent interval $[a,b)$ & $n$ & Wu weight\\
\midrule
$7$ & $p_1$ & $[\sigma_1,1/3)$ & $p_2p_3$ & $2$\\
$8$ & $p_1$ & $[\sigma_2,1/3)$ & $p_2p_3$ & $2$\\
$9{\rm h}$ & $p_1$ & $[1/8,\sigma_1)$ & $p_2p_3$ & $1$\\
$9{\rm l}$ & $p_2$ & $[\sigma_1,11/24)$ & $p_1p_3$ & $1$\\
$10$ & $p_3$ & $[\sigma_1,1-\kappa_2-\sigma_2)$ & $p_1p_2$ & $1$\\
\bottomrule
\end{tabular}
\end{center}
The two parts of row $9$ correspond to $p_1\ge N^{1/8}$ and
$p_1<N^{1/8}$ respectively.  In the latter case
$p_2<\sqrt{N/p_1}\le N^{11/24}$.

For a family with endpoints $a,b$ put
\[
 M_{a,b}=\left\lceil\frac{(b-a)L}{\log(1+\eps_0)}\right\rceil,
 \quad \omega_j=N^a(1+\eps_0)^j,\quad 0\le j<M_{a,b}.
\]
\[
 X_j=N/\omega_j,\quad Z_j=\omega_j,\quad
 Y_j=\min\{(1+\eps_0)\omega_j,N^b\}.
\]
Every selected prime lies in exactly one of these half-open intervals,
including when the displayed quotient is an integer.  The final interval
is clipped, not supplemented by an overlapping endpoint box.
The coefficient $a_j(n)$ is the indicator of the following set, with
$n<X_j$ in every case:
\[
\begin{array}{ll}
7,8: & n=p_2p_3,\quad \omega_j\le p_2\le p_3,\\
9{\rm h}: & n=p_2p_3,\quad N^{\sigma_1}\le p_2\le p_3,\\
9{\rm l}: & n=p_1p_3,\quad z_*\le p_1<N^{1/8},\
                         p_3\ge\omega_j,\ p_1\omega_j^2<N,\\
10: & n=p_1p_2,\quad
 N^{\kappa_2}\le p_1<N^{\sigma_2}\le p_2,\ p_1p_2^2<N.
\end{array}
\]
Factorization recovers the unordered pair of prime factors of $n$.
In the first three cases the increasing order recovers the labels; in
$9{\rm l}$ and $10$ the disjoint prime ranges recover them.
Thus $a_j(n)\in\{0,1\}$, including the case of a prime square.  These
conditions do not depend on the selected prime inside its box.
The original positive source is contained in this box source, since
$pn<N$ implies $\omega_j n<N$.

Define the actual positive integer multiset
\[
 \B_j=\{|N-pn|:n<X_j,\ a_j(n)=1,\ Z_j\le p<Y_j\}.
\]
Here $pn$ is odd and $N$ is even, so $N-pn\ne0$.
The use of the absolute value preserves divisibility by every modulus
and makes the sieve theorem applicable to positive integers.
It also permits the overhang $N<pn<(1+\eps_0)N$ without an unproved
strip-distribution estimate.

For the tag $(p,n)$ use model mass one and
\[
 g_{p,n}(q)=
 \begin{cases}0,&q\mid pn,\\1/(q-1),&q\nmid pn,\end{cases}
 \qquad(q\nmid N).
\]
Then the discrepancy in \cref{lem:model} is precisely the expression
in \eqref{eq:rectangle-discrepancy}, since
$g_{p,n}(d)=\ind_{(pn,d)=1}/\varphi(d)$ for square-free $d$.
There is no substitution of a common density before this identity.
Every prime factor of $pn$ is at least $z_*$ and there are at most three,
so, for $w=y$,
\begin{equation}\label{eq:box-tag-product}
 V_{p,n}(y)\le V_N(y)(1+24/z_*).
\end{equation}
Indeed each deleted factor is $(q-1)/(q-2)\le1+2/z_*$ and
$(1+2/z_*)^3<1+24/z_*$.  The same comparison holds after removal of an
exceptional $q_1<z_*$.

Apply \cref{lem:model} at $D=N^\theta$, $w=y$.
All boxes satisfy
\[
 N^{0.398}\le X_j\le N^{7/8},\qquad Z_j\ge N^{1/8},\qquad
 N< X_jY_j\le(1+\eps_0)N.
\]
By \eqref{eq:level-support}, $QD<D_R$ for every box.
In the large rectangle-exceptional case delete its largest prime $q_1$.
Its threshold exceeds $\prod_{2<p\le31}p$, so $q_1\ge37$ and the
product inflation is at most $36/35$.  No equality between this exceptional conductor and the
one in the ordinary rows is assumed.  The argument $3\theta-\Delta$
is smaller than the actual argument $3\theta$, hence is safe in either
case.

Writing the table's weights as $\lambda_{a,b}$, the complete rectangle
remainder is therefore bounded by
\begin{equation}\label{eq:box-error}
 \frac{E_{\rm rect}}{4\Th}
 \le\frac{(1+\eps_0)m_N}{4(1/2)L}
 \sum_{\rm families}\frac{\lambda_{a,b}}{a^3}
 \left(\frac{(b-a)L}{\log(1+\eps_0)}+1\right)
 <0.0085.
\end{equation}
The evaluated left side is below $0.008032$ at $c=24.9$.
This is a sum over all five families and every box, not a per-box error.
In particular, there is no implicit factor six in the coefficient norm.

We record the main integrals.  Counting the remaining third prime by
\[
 \pi\left(\frac{(1+\eps_0)N}{p_1p_2}\right)
 \le\frac{1.00005(1+\eps_0)N}
              {L p_1p_2(1-t-u)},\qquad
 t=\log p_1/L,\quad u=\log p_2/L,
\]
is legitimate throughout these boxes: $1-t-u>1/4$ and the logarithm
of every prime-counting argument is at least $L/4>30002$.
In particular $(1-6/L)^{-1}<1.00005$, as checked by the certificate.
The changes in the exponent domains are bounded in
\cref{sec:transfer}.  At the unbroadened domain define
\begin{align}
 J_7&=\int_2^{1/\sigma_1-1}\frac{\log(v-1)}v\,dv,&
 J_8&=\int_2^{1/\sigma_2-1}\frac{\log(v-1)}v\,dv,\label{eq:J78}\\
 J_9&=\int_{\kappa_1}^{\sigma_1}
       \frac{\log((1-t-\sigma_1)/\sigma_1)}{t(1-t)}\,dt,&
 J_{10}&=\int_{\kappa_2}^{\sigma_2}
       \frac{\log((1-t-\sigma_2)/\sigma_2)}{t(1-t)}\,dt.\label{eq:J910}
\end{align}
For example, the inner integral in rows $7,8$ is
$\int_t^{(1-t)/2}du/[u(1-t-u)]
 =\log((1-2t)/t)/(1-t)$; the change $v=1/t-1$ gives
\eqref{eq:J78}.  For rows $9,10$ replace the lower inner endpoint by
$\sigma_1,\sigma_2$ respectively, giving \eqref{eq:J910}.
The upper main quantities used in the ledger are
\begin{equation}\label{eq:Bi}
 B_i=\frac{36}{35}\,6e^{-\gamma}\Fup(3\theta-\Delta)J_i
 \qquad(i=7,8,9,10).
\end{equation}
The factor $1.00005(1+\eps_0)$, local products, and finite deletions have
not been hidden in these integrals; they are charged separately below.

\section{Finite transfers and local terms}\label{sec:transfer}

\subsection{A prime-measure estimate with explicit variation}

Let $\mu_N$ be the measure on $[\kappa_1,1]$ putting mass $1/p$ at
$\log p/L$ for $z_*\le p<N$.  Let $d\mu(t)=dt/t$ there.
BJS's reciprocal-prime inequalities, including the at-most-one-prime
endpoint correction, give
\begin{equation}\label{eq:measure-discrepancy}
 \sup_{\kappa_1\le t\le1}
 \left|(\mu_N-\mu)([\kappa_1,t))\right|
 \le d_N:=\frac{3\cdot10^{-6}}{\kappa_1L}+2N^{-\kappa_1},
 \qquad \mu_N([\kappa_1,1])<3.
\end{equation}
The constant $3\cdot10^{-6}$ bounds
$2.964\cdot10^{-6}+1.436\cdot10^{-16}$, not twice the sum:
the two endpoint errors each lie in the same asymmetric interval,
so their difference has at most that interval's length.

Here is the finite-dimensional transfer we use.  Suppose the coordinate
sections of a domain are unions of at most ten intervals, and its kernel
$G$ has magnitude at most $M$, at most ten additional jump points in
each section, and derivative at most $M'$ away from those points.
Integration by parts against the cumulative discrepancy in
\eqref{eq:measure-discrepancy} bounds a one-coordinate difference by
$d_N(40M+M')$.  The coefficient $40M$ includes both domain endpoints
and both sides of every jump.  Telescoping the product measures gives
\begin{equation}\label{eq:transfer}
 \left|\int G\,d\mu_N^{\,d}-\int G\,d\mu^{\,d}\right|
 \le d\,3^{d-1}d_N(40M+M').
\end{equation}
To see the telescoping explicitly, replace coordinates $1,2,\ldots,d$
one at a time.  Condition on the other $d-1$ coordinates, apply the
one-dimensional integration-by-parts bound, and integrate that bound
against their product measure, whose mass is at most $3^{d-1}$.
This proof also covers strict inequalities and diagonal atoms.

The following bounds apply to the kernel excluding the factors $dt/t$.
They include the displayed sieve prefactors, but not the Wu row weight.
\begin{center}
\begin{tabular}{@{}lrrr@{}}
\toprule
Kernel & $d$ & $M$ & $M'$\\
\midrule
$A_1$ & 1 & 120 & 20000\\
$A_2,A_3$ (each) & 1 & 120 & 2000\\
$L_4,L_5$ (each) & 2 & 48 & 3000\\
Each triple-prime box mass & 2 & 128 & 512\\
Each switched four-prime mass & 4 & 8000 & 0\\
\bottomrule
\end{tabular}
\end{center}
These are analytic bounds, not values sampled on a mesh.
On the domains at issue, $1/t\le12$, $1/t^2\le144$,
$F\le4$, $f\le1$, $|F'|,|f'|\le5$, and $|h'|\le1$.
In fact $\Fup\le4$ and $|\Fup'|\le6$: the first follows from
$2e^\gamma+106\eta<4$, and the second from
$106\eta e^2<1$.  These inequalities are included in the certificate.
The delay equations prove the derivative bounds directly, using
$s\ge1$ (and $s-1\ge1$ whenever its nonzero term is needed).
For $\Flow$, $|\Flow|\le2$ and $|\Flow'|\le10$ on its smooth pieces.
The arguments in $A_1$ have derivative at most $72$, and the other
ordinary arguments at most $12$; the prefactor
$2e^{-\gamma}/\kappa_1<14$ then gives the first three lines.
In a box kernel, $1-t-u>1/4$ and
$(36/35)6e^{-\gamma}\Fup(3\theta-\Delta)<16$, giving
$M<64$, $M'<256$; the table doubles these to cover the inflated box.
The switched multiplier is less than $8000$ using
$2e^{-2\gamma}/\kappa_1^2<100$ and $F\le4$.
All boundaries are linear exponent inequalities; their sections are
intervals.  Breaks of the sieve functions at $2,3,4,5,6$ and the jump
of $\Flow$ at $2$ use fewer than ten additional points.
These facts prove every entry of the table.

The sum of the right sides of \eqref{eq:transfer}, including the Wu
weights and the two switched rows, is less than $10^9d_N$ even before
division by four.  We allow $10^{10}d_N$ in the normalized ledger.

\subsection{Box overhangs, prime counting and products}

Put $h_0=\log(1+\eps_0)/L$.  For rows $7,8$, the box condition gives
$p_2\ge p_1/(1+\eps_0)$ and $p_1p_2^2<(1+\eps_0)N$.
Thus the exponent boundaries move from $u=t$ to $u\ge t-h_0$ and
from $t+2u<1$ to $t+2u<1+h_0$.
For $9{\rm h}$ the lower condition $u\ge\sigma_1$ is unchanged and
only the latter boundary moves.  For $9{\rm l}$,
$p_3\ge p_2/(1+\eps_0)$ and
$p_1p_2^2<(1+\eps_0)^2N$; the upper boundary moves by at most $2h_0$.
For $10$, $p_1p_2^2<N$ is retained exactly and the product bound only
enlarges to $p_1p_2p_3<(1+\eps_0)N$.
The $p_1=N^{1/8}$ split is disjoint and makes no additional mass.
The five families consequently require at most eight strips, each of
coordinate width at most $2h_0$.  Their continuous mass is bounded using
$dt/t\le12\,dt$ and the table.  Their discrete mass first uses
\eqref{eq:transfer}, already budgeted above.  A common conservative
bound for all these continuous strips is $10^{10}\cdot16h_0$.

The raw total of all negative main terms in the final ledger is below
$60$, certified with Arb.  Since
$1.00005(1+\eps_0)<1.0002$, the common prime-count inflation costs at
most $60/(4\cdot5000)=0.003$.  Therefore
\begin{equation}\label{eq:transfer-cost}
 E_{\rm transfer}/(4\Th)
 \le10^{10}(d_N+16h_0)+0.003<0.0035.
\end{equation}
This also permits the factor $1.00005$ in ordinary upper model masses.
The evaluated proxy is less than $0.003251$ at $c=24.9$.

The product calculation in BJS Lemma 38, or its displayed factorization
into the full Mertens product and the local factors of $N$, gives
\begin{equation}\label{eq:Vlocal}
 V_N(N^t)=\frac{2e^{-\gamma}C_N}{tL}(1+\epsilon_t),
 \quad |\epsilon_t|\le
 \frac{2}{(\kappa_1L)^2}+10L N^{-\kappa_1},
 \qquad t\ge\kappa_1.
\end{equation}
To specify the tail used here: the full prime product contributes at
most $1/\log^2(N^t)$; the tail of
$\prod_{p>2}(1-1/(p-1)^2)$ is bounded by
$2/(N^t-2)$; and the omitted local factors with $p\mid N$, $p\ge N^t$
are bounded in logarithm by $2\omega(N)/(N^t-2)$.
Using $\omega(N)\le L/\log2$ and $e^x\le1+2x$ for $0\le x\le1/2$
gives \eqref{eq:Vlocal}.  The same bound covers all prefix products,
with $\log(Nm_j)\le L+10c$; their local ratios are already retained
exactly in \cref{lem:exceptional}.
The full integer-sieve product satisfies
\begin{equation}\label{eq:Vfull}
 \prod_{p<z_*}(1-1/p)\le
 (1+\eta)e^{-\gamma}/\log z_*.
\end{equation}
For example BJS's quoted Mertens bound
$e^{-\gamma}(1+\log^{-2}z_*)/\log z_*$
is stronger, since $(\kappa_1L)^{-2}<\eta$.
There is no $N$-dependent local factor in \eqref{eq:Vfull}.

For a tag of one or two primes at least $z_*$,
$a/\varphi(a)\le(1-1/z_*)^{-2}<1+4/z_*$.
Together with \eqref{eq:box-tag-product}, \eqref{eq:Vlocal}, and the
total kernel bounds above, all such normalized local corrections are
bounded by
\begin{equation}\label{eq:local-cost}
 E_{\rm local}/(4\Th)
 \le1000\left(\frac{2}{(\kappa_1L)^2}
                  +40L N^{-\kappa_1}+L^{-9}\right)<10^{-6}.
\end{equation}
The $L^{-9}$ allowance here is harmless slack, not a second charge of
the integer rough-sieve remainder, which is assigned to the switch.
In lower rows where the chosen lower function is negative, use $S\ge0$
before making any mass or product replacement.

\subsection{A single finite combinatorial envelope}

All combinatorial exclusions used here are bounded, before division
by $4\Th$, by
\begin{equation}\label{eq:finite-envelope}
\begin{split}
 E_{\rm finite}:={}&(4+7W^2)B+6W^2N^{2/3}+24N^{1/3}\\
 &+4(1+W)^5N^{\kappa_1}
              +100(1+W)^6N^{1-\kappa_1},
 \qquad W=L/\log2.
\end{split}
\end{equation}
Here are the assignments, so no unspecified terminal is needed.
The first $(4+W^2)B$ is \cref{lem:wu}.  When rows $7$--$10$ are
converted to three-prime sources, the removed nonsquare-free source
members contribute at most $6W^2B$: each tag uses two distinct divisors
and the four Wu weights sum to six.  A residual equal to one gives
$p_1p_2<N^{2/3}$ in all four rows, so at most $6W^2N^{2/3}$ is lost.
For a complementary prime below $y$, there are at most $y$ source
values and at most four ordered triple tags per value; after the
weights this is at most $24y$.

In rows $13,14$, a complementary prime below $z_*$ has at most
$\binom{\omega(a)}4\le(1+W)^4$ ordered tags for its source
$a=N-p$.  Allowing also the at most $W$ complementary primes dividing
$N$ gives less than $4(1+W)^5z_*$.  The latter allowance is unnecessary
for our definition of $\A$, but explicitly covers the more general
local implication in \cref{lem:switch-local}.
Finally, removing a prime $p\mid N$, $p\ge z_*$ from a lower prime
integral costs a reciprocal mass at most $W/z_*$.  Multiplying by
the kernel bounds, the other coordinate masses, $N/L^2$, and all
Wu weights is less than the last term in \eqref{eq:finite-envelope}.
That term also covers any endpoint prime or repeated-prime boundary
that is removed before using an ordered source.  Ordered-source
diagonals in the integral transfers are already allowed by
\eqref{eq:transfer}, so this is conservative.
All floor contributions in these counts use $N/q^2+1$, as in $B$;
they have not been replaced by $N/q^2$.
Since $C_N>1/2$, direct comparison of the five displayed terms gives
\begin{equation}\label{eq:finite-cost}
 E_{\rm finite}/(4\Th)\le e^{30}L^{12}N^{-1/12}<10^{-6}.
\end{equation}
The explicit logarithmic comparison and its half-line validity are
included in the certificate and \cref{sec:uniform}.

\section{The complementary-prime switch}\label{sec:switch}

\subsection{Local implication and unique factor recovery}

Let $\mathcal C_{13},\mathcal C_{14}$ be the multisets of tuples
$(p_1,p_2,p_3,p_4,r_0)$ satisfying the respective prime ranges in
$U_{13},U_{14}$, and
\[
 r_0\ge1,\quad (r_0,P_{\rm all}(p_2))=1,\qquad
 p_1p_2p_3p_4r_0\le N-2.
\]
Do not impose coprimality with $N$ in these enlarged sets.
Set $b=N-p_1p_2p_3p_4r_0$, so $b$ is a positive odd integer.
Tags are retained even if their values $b$ coincide.

\begin{lemma}[The local switch, with its exceptional alternative]
\label{lem:switch-local}
Write $\varpi=N-p_1p_2p_3p_4r_0$ and suppose the four primes are
coprime to $N$.  If $\varpi$ is prime and
$(r_0,P_N(p_2))=1$, then either $\varpi\mid N$ or
$(r_0,N)=1$ and $(r_0,P_{\rm all}(p_2))=1$.
For $\varpi\ge z_*$ in the latter case, the complementary upper
Rosser weight at $z_*$ is at least one.
\end{lemma}
\begin{proof}
If a prime $q$ divides $(r_0,N)$, it divides $\varpi$, hence
$q=\varpi\mid N$.  Otherwise $(r_0,N)=1$; every prime below $p_2$
dividing $N$ is then absent from $r_0$, and every other such prime is
absent by the original sift.  This is precisely full $p_2$-roughness.
The upper Rosser divisor sum bounds the sifted indicator pointwise;
a prime $\varpi\ge z_*$ has no prime factor below $z_*$.
Its divisor sum is therefore at least one, and is nonnegative on
every positive integer.  This last fact permits enlargement to
$\mathcal C_i$ after the implication has been applied.
\end{proof}

For $\A$ itself $\varpi\nmid N$, so only the main alternative occurs.
Apart from the finite small-prime charge in \eqref{eq:finite-envelope},
$U_i$ is bounded by the sum of the upper Rosser weights on
$\mathcal C_i$.  We retain full $p_2$-roughness throughout the
distribution proof.  It is weakened to $z_*$-roughness only when
bounding the principal model mass.

Regroup the source by
\[
 n=p_3p_4,\qquad m=p_1p_2r_0.
\]
The two prime factors of $n$ determine $(p_3,p_4)$, since $p_3<p_4$.
The two least prime factors of $m$, counted with multiplicity, determine
$(p_1,p_2)$: $p_1<p_2$ and every prime factor of $r_0$ is at least
$p_2$.  Dividing $m$ by $p_1p_2$ then recovers $r_0$.
Thus both coefficient sequences are exactly bounded by one.
This remains true when $p_2\mid r_0$ or $r_0=1$; no square-free
condition on $r_0$ is needed.  The only cross-conditions between the
two groups are
\begin{equation}\label{eq:switch-cross}
 p_2<p_3,\qquad mn\le N-2.
\end{equation}
All other prime ranges, including $p_3p_4<N^{\lambda}$ in row $14$,
belong to one group.  In particular
$N^{1/6}\le n<N^{\lambda}<N^{1/3}$.

Here are the coefficient formulas that will be used below.  Let
\begin{align*}
 \mathscr N_{13}&=\{(p_3,p_4):z_*\le p_3<p_4<N^{\kappa_2}\},\\
 \mathscr N_{14}&=\{(p_3,p_4):z_*\le p_3<N^{\kappa_2}\le p_4,
                                  \ p_3p_4<N^\lambda\},\\
 \mathscr M&=\{(p_1,p_2,r_0):z_*\le p_1<p_2<N^{\kappa_2},\ r_0\ge1,
                         \ (r_0,P_{\rm all}(p_2))=1\}.
\end{align*}
For a dyadic block $I_X=[X,2X)$, clipped at $N^{1/3}$, put
$M_X=N/X$ and define
\begin{align*}
 \alpha_{i,X}(n)&=
  \#\{(p_3,p_4)\in\mathscr N_i:p_3p_4=n,\ n\in I_X\},\\
 \beta_X(m)&=
  \#\{(p_1,p_2,r_0)\in\mathscr M:p_1p_2r_0=m,\ m\le M_X\}.
\end{align*}
The unique recovery just proved says, pointwise,
\[
 \alpha_{i,X}(n),\beta_X(m)\in\{0,1\},\qquad
 \sum_n|\alpha_{i,X}(n)|^2<2X,\quad
 \sum_m|\beta_X(m)|^2\le M_X.
\]
On their supports denote the recovered primes by $p_3(n)$ and $p_2(m)$.
Because $n\ge X$, a pair with $m>M_X$ cannot pass $mn\le N-2$;
therefore the cutoff in $\beta_X$ loses nothing.  For every character
$\chi$, the following is an exact finite identity:
\begin{equation}\label{eq:switch-source-identity}
 \sum_{\mathcal C_i}\chi(mn)
 =\sum_X\sum_{n,m}\alpha_{i,X}(n)\beta_X(m)\chi(nm)
   \ind_{p_2(m)<p_3(n)}\ind_{mn\le N-2}.
\end{equation}
The blocks are disjoint, so this identity contains neither an ordering
multiplicity nor a hidden convolution energy estimate.

\subsection{Exact character decomposition and the low conductors}

Use $D_\circ=N^{2/5}$, $w=z_*$ and $D_*=QD_\circ$.
For the switched exceptional alternative use the conductor cutoff
\[
 X_q=z_*/2,\quad w_q=\log x_2(X_q),\quad
 q_0=\lceil w_q^{10}\rceil-1.
\]
Thus $w_q^{10}/2\le q_0<w_q^{10}$, including at integral values
of $w_q^{10}$.  If the relevant conductor up to $q_0$ is at least
$w_q^\delta$, remove its largest prime $q_1$ from the sieve prime set.
The threshold exceeds $\prod_{2<p\le31}p$ and hence its largest prime
is at least $37$; its product cost is at most $36/35$.  In the other
case make no deletion.
The zero bound $\nu$ from \eqref{eq:vformula} is valid in both cases.

\begin{theorem}[Finite switched distribution]\label{thm:switch}
For the two explicitly defined tuple multisets $\mathcal C_i$, use
all square-free $d<D_*$ with $d\mid P_N(z_*)$ in the small switched
alternative, and omit multiples of $q_1$ in the large alternative.
For the discrepancies $r_i(d)$ defined in \eqref{eq:switch-character},
\[
 \sum_d\bigl(|r_{13}(d)|+|r_{14}(d)|\bigr)
 \le E_{\rm low}+E_{\rm high}+E_{\rm P},
\]
with the complete formulas \eqref{eq:switch-low},
\eqref{eq:switch-high} and \eqref{eq:switch-perron}.
Together with \eqref{eq:rough-count} this proves the switched upper
main term \eqref{eq:Bswitch}, with only the transfers, local products
and finite exclusions explicitly assigned in \cref{sec:transfer}.
\end{theorem}
The proof follows; in particular, the theorem does not assume a
weighted bilinear distribution result.

Every source factor of $mn$ is at least $z_*$; every prime factor of a
supported modulus is below $z_*$.  Consequently $(mn,d)=1$
identically on the support.  Character orthogonality gives the exact
remainder
\begin{equation}\label{eq:switch-character}
\begin{split}
 r_i(d)&=\#\{(m,n)\in\mathcal C_i:mn\equiv N\pmod d\}
                         -|\mathcal C_i|/\varphi(d)\\
 &=\frac1{\varphi(d)}\sum_{\chi\ne\chi_0}\overline{\chi}(N)
                    \sum_{(m,n)\in\mathcal C_i}\chi(mn).
\end{split}
\end{equation}
For square-free $d=rs$, $(r,s)=1$, each character has exactly one
primitive inducing character modulo $r$.  Since $(mn,s)=1$,
induction introduces no correction on this source.  In particular
the often missing term $p_4\mid s$ is identically zero, not merely
small.  Therefore
\begin{equation}\label{eq:switch-primitive}
 \sum_{d<D_*}|r_i(d)|
 \le\sum_{rs<D_*}\frac{\mu^2(rs)}{\varphi(r)\varphi(s)}
       \sum_{\substack{\chi\bmod r\ {\rm primitive}\\\chi\ne\chi_0}}
       \left|\sum_{\mathcal C_i}\chi(p_1p_2p_3p_4r_0)\right|.
\end{equation}
The modulus set retains $(d,N)=1$ and the exceptional deletion.
We drop such restrictions only inside nonnegative upper bounds.

For $r\le q_0$, fix $p_1,p_2,p_3,r_0$ first.  The allowed $p_4$
form a single interval with lower endpoint at least $z_*$ and upper
endpoint at most $N/(p_1p_2p_3r_0)$.
For row $13$ intersect $(p_3,N^{\kappa_2})$ with that upper bound;
for row $14$ intersect $[N^{\kappa_2},N^{\lambda}/p_3)$ with it.
Empty intervals contribute zero.  Both prime-counting endpoints,
rounded to the corresponding integers, exceed $X_q$.
BJS Lemma 29 applies before any Perron integral is introduced.
Its bound at the two endpoints, and $\log p_4\ge\kappa_1L$, give
at most
\[
 \frac{2v(X_q)N}{(\kappa_1L)^5p_1p_2p_3r_0}
\]
for the sum of absolute primitive-character prime sums.
The reciprocal-prime mass of $[N^{\kappa_1},N^{\kappa_2})$ is
less than $1/3$ by \eqref{eq:measure-discrepancy}.  The ordered
three-prime reciprocal sum is at most its cube divided by $6$,
and is therefore below one.  Also
$\sum_{r_0\le N}1/r_0\le1+L$.  Finally
$\sum_{r\le q_0}\mu^2(r)/\varphi(r)\le1.1\log(w_q^{10})$ and
$\sum_{s<D_*}\mu^2(s)/\varphi(s)\le1.1\log D_*$.
Thus the two rows together satisfy
\begin{equation}\label{eq:switch-low}
 E_{\rm low}\le
 \frac{4v(X_q)N(1+L)}{(\kappa_1L)^5}
 \,1.21\log(w_q^{10})\log D_*.
\end{equation}
This is a bound for all low conductors and both variable endpoints.
At $c=24.9$ it is below $1.076\cdot10^{-12}\Th$ using the
uniform envelope $v(X_q)\le40/L$.

\subsection{Perron separation, with an explicit tail}

Divide $n$ into intervals $X\le n<2X$, beginning at $N^{1/6}$ and
clipping the last at $N^{1/3}$.  There are at most
$K_n=2+L/(6\log2)$ intervals.  On one interval write $M=M_X=N/X$
and abbreviate the coefficients in \eqref{eq:switch-source-identity}
as $\alpha(n),\beta(m)$.
The two exact tests in \eqref{eq:switch-cross} are
\[
 \ind_{\,p_3/(p_2+1/2)>1},\qquad
 \ind_{\,(N-3/2)/(mn)>1}.
\]
They avoid equality at the integration boundary for every integer
tuple.  Put $\sigma=1/L$, $T=N^4$.  For $x>0$, $x\ne1$, let
\[
 \mathscr P_T(x)=\frac1{2\pi i}
                \int_{\sigma-iT}^{\sigma+iT}\frac{x^s}{s}\,ds.
\]
Contour integration to the right or left of the pole at zero, followed
by integration by parts on the two horizontal rays, gives
\begin{equation}\label{eq:Perron-tail}
 |\mathscr P_T(x)-\ind_{x>1}|
 \le\frac{4x^\sigma}{T|\log x|}.
\end{equation}
For precision, the untruncated vertical integral equals the indicated
residue; for a tail ray
$\int_T^\infty e^{it\log x}/(\sigma+it)\,dt$, integration by parts
bounds its absolute value by $2/(T|\log x|)$.  The two rays and
$1/(2\pi)$ give the weaker constant $4$ above.

Both ratios lie between $1/(2N)$ and $2N$, have
$|\log x|\ge1/(4N)$, and $x^\sigma\le e^2$.
The latter two facts follow from their half-integer gaps and all
participating integers being at most $2N$.
Their individual tails are at most $16e^2N/T$.
As $16e^2N/T<1$, replacing the product of the two indicators by the
product of the two integrals costs at most $64e^2N/T$ per pair.
There are at most $2XM=2N$ pairs per block, fewer than $D_*$
moduli, and two rows; after summing the character weights
$\varphi(d)^{-1}\sum_\chi1\le1$, the full tail is at most
\begin{equation}\label{eq:switch-perron}
 E_{\rm P}\le256e^2K_nN^2D_*/T.
\end{equation}
It is bounded after normalization by the even larger envelope
$e^{20}L^{20}D_*N^{-2}$ in the certificate.

The two remaining integrals separate as a product of a sum in $n$
and a sum in $m$.  Their weights are
$p_3^s n^{-t}$ and $(p_2+1/2)^{-s}m^{-t}$, with the scalar
$(N-3/2)^t$ outside.  Here $\Re s=\Re t=\sigma$.
Each weight is bounded by $e^2$, and the outside scalar by $e$.
Consequently the square norms and the outside factor together cost
less than $e^6$, uniformly in both imaginary parts.
The total variation of each Perron kernel is
\[
 \frac1{2\pi}\int_{-T}^T\frac{dt}{|\sigma+it|}
 =\frac1\pi\operatorname{arsinh}(T/\sigma)<2L.
\]
The certificate verifies the endpoint inequality; its difference is
increasing on the half-line.  No maximal-sum theorem or complex
small-conductor prime estimate is being assumed.

\subsection{The four large-conductor terms are all added}

For $r>q_0$ split into $R\le r<2R$ and use the multiplicative large
sieve.  Zero-pad the $n$-sequence to the prefix $n<2X$.  Then
\[
 (2X+12R^2)\sum_n|\alpha(n)|^2
 <(2X+12R^2)2X\le4(X+12R^2)X,
\]
whereas the $m$-sequence contributes at most
$(M+12R^2)M$.  Thus the square root costs the explicit factor $2$
in the next display; it also absorbs every endpoint $+1$.
Cauchy--Schwarz, $|\alpha|,|\beta|\le1$, and the weights just bounded
therefore give the following majorant before the two Perron factors:
\[
 2e^6\frac{\{(X+12R^2)(M+12R^2)XM\}^{1/2}}R.
\]
In particular, the floor $+1$ is already covered by the norm
$\sum_{n<2X}|\alpha(n)|^2<2X$ and
$\sum_{m\le M}|\beta(m)|^2\le M$.
The elementary inequality
\[
 \frac{\{(X+12R^2)(M+12R^2)XM\}^{1/2}}R
 \le\frac{XM}R+\sqrt{12X}\,M+\sqrt{12M}\,X+12R\sqrt{XM}
\]
has four positive terms.  Summing $1/R$ geometrically gives at most
$2/q_0$; summing $R$ gives at most $2D_*$; the number of intervals
is at most $K_R=2+0.4154L/\log2$.  Define
\[
 B_1=2XM/q_0,\quad B_2=K_R\sqrt{12X}\,M,\quad
 B_3=K_R\sqrt{12M}\,X,\quad B_4=24D_*\sqrt{XM}.
\]
After both rows, all $n$ blocks, both Perron factors, and the
$s$-sum in \eqref{eq:switch-primitive}, a conservative bound is
\begin{equation}\label{eq:switch-high}
 E_{\rm high}\le64e^6K_n(2L)^2(1.1\log D_*)
                 (B_1+B_2+B_3+B_4).
\end{equation}
In this expression $XM=N$, $X\ge N^{1/6}$, $M\ge N^{2/3}$,
$q_0\ge w_q^{10}/2$, and $D_*<N^{0.4154}$.
These substitutions give a quantity below $10^{-20}\Th$ for every
$c\ge24.9$; at the endpoint the certificate encloses it below
$1.597\cdot10^{-28}\Th$.  Induced characters have been counted
once through $d=rs$, and there is no uncharged ordering multiplicity.

\subsection{The integer rough sieve and the switched main mass}

Only now replace the condition on $r_0$ by full $z_*$-roughness.
For $P=p_1p_2p_3p_4$ let $R=N/P$, enlarging the upper endpoint
$(N-2)/P$ when necessary.  In both rows
\[
 R\ge N^a,\qquad
 a=1-\lambda-2\kappa_2
  =\frac12+2\kappa_1-2\kappa_2=\frac{999}{2000}.
\]
For row $13$ this follows from $p_1p_2p_3p_4<N^{4\kappa_2}$,
which is stronger.  For row $14$ it follows directly from
$p_3p_4<N^{\lambda}$ and $p_1p_2<N^{2\kappa_2}$.
Use the integer sequence $1\le r_0\le R$, model total $R$,
density $g(q)=1/q$, and
\[
 D_r=\frac{R}{Q_{\rm all}\log^{10}R},\qquad
 Q_{\rm all}=\prod_{p\le10^9}p.
\]
For every $d$, the discrepancy is $\lfloor R/d\rfloor-R/d$ and its
absolute value is at most one.  There are fewer than
$Q_{\rm all}D_r=R/\log^{10}R$ supported positive integers $d$.
Moreover
\[
 \frac{\log D_r}{\log z_*}
 \ge\frac{aL-\overline q-10\log(aL)}{\kappa_1L}
 >s_r:=\frac{12037}{2500}.
\]
The monotonicity of $x-10\log x$ for $x>10$ proves the first
inequality.  The second holds uniformly by \cref{sec:uniform}.
The model sieve and \eqref{eq:Vfull} give
\begin{equation}\label{eq:rough-count}
 \#\{r_0\le R:(r_0,P_{\rm all}(z_*))=1\}
 \le\frac{(1+\eta)e^{-\gamma}}{\kappa_1L}\Fup(s_r)R
                       +\frac{R}{\log^{10}R}.
\end{equation}
This is a statement about the full prime product, not $P_N(z_*)$.

The complementary model is $|\mathcal C_i|$ with common density
$1/(q-1)$.  The upper model sieve at $D_\circ$ has argument
$\log D_\circ/\log z_*=800/191$.
To use a single conservative expression in both exceptional
alternatives, replace it by $800/191-\Delta$ and include $36/35$.
Summing the main term of \eqref{eq:rough-count} over the prime ranges
and applying the already budgeted transfer gives
\begin{equation}\label{eq:Bswitch}
 B_{\rm sw}=\frac{36}{35}(1+\eta)
       \frac{2e^{-2\gamma}}{\kappa_1^2}
       \Fup(s_r)\Fup(800/191-\Delta)(I_{13}+I_{14}),
\end{equation}
where
\[
 I_{13}=\frac{\log^4(\kappa_2/\kappa_1)}{24},\qquad
 I_{14}=\int_{\kappa_1}^{\kappa_2}
       \frac{\log^2(t/\kappa_1)}{2t}
       \log\left(\frac{\lambda-t}{\kappa_2}\right)\,dt.
\]
Indeed the first is the ordered four-simplex of the measure $dt/t$.
In the second, fixing $t_3=t$ gives
$\int_{\kappa_1<t_1<t_2<t}dt_1dt_2/(t_1t_2)
=\tfrac12\log^2(t/\kappa_1)$, and then integrate
$\kappa_2\le t_4<\lambda-t$.

The remainder of \eqref{eq:rough-count}, after the complementary
sieve multiplier, is at most $e^{20}L^{-9}\Th$.
For a direct bound, its reciprocal four-prime sum is below
$3^4=81$, $\log R\ge aL$, $V_N(z_*)\le14C_N/L$,
and $(36/35)\Fup(800/191-\Delta)<3$.
These yield $3402a^{-10}L^{-9}\Th<e^{20}L^{-9}\Th$.
All constants here are checked as rational/Arb inequalities.

Combining \eqref{eq:switch-low}, \eqref{eq:switch-perron},
\eqref{eq:switch-high}, and this last rough remainder proves the
required complete distribution and switching estimate.  Their sum,
even without using the final division by four, is below
$10^{-6}\Th$.  The ledger uses the allowance
\begin{equation}\label{eq:switch-cost}
 E_{\rm sw}/(4\Th)<10^{-6}.
\end{equation}
Small complementary primes and local source exclusions remain assigned
to \eqref{eq:finite-envelope}; the switch error is not counted twice.

\section{One forward ledger and the assembly}\label{sec:assembly}

All the numbers in this section are enclosed by the Arb computation
\cert{c249_verify.py}; no decimal is fitted to a desired final answer.
The symbols in the first column are exactly those in
\eqref{eq:L0}--\eqref{eq:Bi} and \eqref{eq:Bswitch}.
The following intervals are rounded outwards to twelve decimal places.
The signed weight is applied before the final division by four.
\begin{center}
\small
\begin{tabular}{@{}lrrc@{}}
\toprule
Quantity & Lower endpoint & Upper endpoint & Signed Wu weight\\
\midrule
% ledger:L0:11.738503274673:11.738503274674:4
$L_0$ & $11.738503274673$ & $11.738503274674$ & $4$\\
% ledger:A1:0.031268140768:0.031268140769:-36/35
$A_1$ & $0.031268140768$ & $0.031268140769$ & $-36/35$\\
% ledger:A2:16.368118871827:16.368118871828:-36/35
$A_2$ & $16.368118871827$ & $16.368118871828$ & $-36/35$\\
% ledger:A3:16.357231429105:16.357231429106:-36/35
$A_3$ & $16.357231429105$ & $16.357231429106$ & $-36/35$\\
% ledger:L4:0.000032859562:0.000032859563:1
$L_4$ & $0.000032859562$ & $0.000032859563$ & $1$\\
% ledger:L5:0.011400175493:0.011400175494:1
$L_5$ & $0.011400175493$ & $0.011400175494$ & $1$\\
% ledger:B7:0.164463777759:0.164463777760:-2
$B_7$ & $0.164463777759$ & $0.164463777760$ & $-2$\\
% ledger:B8:0.166078002505:0.166078002506:-2
$B_8$ & $0.166078002505$ & $0.166078002506$ & $-2$\\
% ledger:B9:6.263354678712:6.263354678713:-1
$B_9$ & $6.263354678712$ & $6.263354678713$ & $-1$\\
% ledger:B10:6.254338521550:6.254338521551:-1
$B_{10}$ & $6.254338521550$ & $6.254338521551$ & $-1$\\
% ledger:Bsw:0.000000172621:0.000000172622:-1
$B_{\rm sw}$ & $0.000000172621$ & $0.000000172622$ & $-1$\\
\bottomrule
\end{tabular}
\end{center}
The box quantities and $B_{\rm sw}$ already include $36/35$; applying
that factor to them a second time would be incorrect.
The three $A_i$ do not include it.  The single master expression is
\begin{equation}\label{eq:master}
 \mathfrak M=
 \frac{4L_0-(36/35)(A_1+A_2+A_3)+L_4+L_5
                  -2B_7-2B_8-B_9-B_{10}-B_{\rm sw}}4.
\end{equation}
Positive rows use lower interval endpoints and negative rows use upper
ones.  Even using only the printed twelve-place intervals gives
$\mathfrak M>0.0235368436443>0.0228$.
Using the full balls gives
\[
 \mathfrak M\in
[\,0.0235368436464714387690418419462,\,
    0.0235368436464714387690418419466\,].
\]
The auxiliary small-exceptional computation is not needed for the
theorem.  Our universal worst-case expression is used even if the
ordinary, rectangular and switched exceptional choices are different.

The following are \emph{allocated upper bounds}, already normalized by
$4\Th$.  Each actual error was bounded before this allocation.
\begin{center}
\begin{tabularx}{\textwidth}{@{}Xlr@{}}
\toprule
Error family & Defining bound & Allocated upper bound\\
\midrule
All rectangle boxes & \eqref{eq:box-error} & $0.008500$\\
Prime measures, box boundaries, prime-count inflation &
 \eqref{eq:transfer-cost} & $0.003500$\\
Ordinary AP remainders & \eqref{eq:ordinary-AP-cost} & $0.000001$\\
Exceptional full-conductor mass & \eqref{eq:exc-cost} & $0.000001$\\
Switch distribution, Perron, integer rough floor &
 \eqref{eq:switch-cost} & $0.000001$\\
Finite combinatorial exclusions & \eqref{eq:finite-cost} & $0.000001$\\
Local sieve products & \eqref{eq:local-cost} & $0.000001$\\
\midrule
Total & & $0.012005$\\
\bottomrule
\end{tabularx}
\end{center}
There is one charge per error family.  In particular the integer rough
floor appears in the switched error and the square-divisor source
term appears in the finite error; neither is subtracted again.

Applying \cref{lem:wu}, the ordinary model estimates, the five box
estimates and the switched estimate, followed by their transfers,
gives for every even $N$ in the range where the seven bounds hold
\begin{equation}\label{eq:assembled}
 \frac{D_{1,2}(N)}{\Th}\ge
 \mathfrak M-\frac{E_{\rm rect}+E_{\rm transfer}+E_{\rm AP}
 +E_{\rm exc}+E_{\rm sw}+E_{\rm finite}+E_{\rm local}}{4\Th}.
\end{equation}
The errors here denote their assigned total upper envelopes, so an
unused allowance is still subtracted.  The next section proves that
every envelope is valid on the full half-line.  It follows that the
right side of \eqref{eq:assembled} is greater than
$0.0228-0.012005=0.010795>0$, proving \cref{thm:main}.
This internal positive margin is not an attempt to optimize an explicit
coefficient.

\section{Uniformity on the entire half-line}\label{sec:uniform}

In this section $c$ varies continuously over $[24.9,\infty)$ and
$L=e^c$.  Establishing bounds on this larger continuous range suffices
for every even integer $N$.  The certificate checks only the starting
inequalities used below; all propagation is analytic.

\subsection{Fixed levels and exceptional alternatives}

The function
\[
 \frac12-e^{-c}(\overline q+20c)
\]
has derivative $e^{-c}(\overline q+20c-20)>0$.
The exceptional level loss $20ce^{-c}/\kappa_1$ has negative derivative.
For $a=999/2000$, the derivative of
\[
 \frac{a-e^{-c}\{\overline q+10\log(ae^c)\}}{\kappa_1}
\]
is $e^{-c}\{\overline q+10\log(ae^c)-10\}/\kappa_1>0$.
These facts prove all uses of $\theta,\Delta,s_r$ on the half-line.
Also $\log D_*/L\le2/5+\overline q/L<0.4154$ everywhere.
The functions $\log x_1(N)=L-5c$ and
$\log x_2(N^b/2)=bL-\log2-15\log(bL-\log2)$, for
$b=1/8,\kappa_1$, are increasing and positive.
Their powers with exponent $\delta$ are therefore increasing.
The endpoint primorial inequalities thus guarantee $q_1\ge37$ for
every large-exceptional branch, irrespective of changes of conductor.
No assertion of monotonicity of the conductor itself is made.

Both exceptional alternatives are covered by the common
$\Flow$ and $36/35$ envelopes; no small-case coefficient is used
to bound a later large case.  Small primes dividing $N$ can change
arbitrarily, but $\log Q\le\overline q$ and the local product comparisons
hold separately for each $N$.  Hence neither changes of exceptional
conductor nor primorial transition points require an unexamined
continuity assertion.

\subsection{A logarithmic derivative bound for the zero terms}

Write $t=bL-\log2$ and $w=t-15\log t$, for
$b=1/8,\kappa_1$.
For every $c\ge24.9$,
\begin{equation}\label{eq:wderivative}
 1<\frac{w'}w<1.01.
\end{equation}
Indeed
\[
 w'-w=\log2+15\log t-15-\frac{15\log2}{t}>0,
\]
and, since $w>t/2$,
\[
 0<\frac{w'}w-1
 <\frac{32\log t}{t}<0.01.
\]
The last bound is checked at the smaller starting $t$;
$\log t/t$ decreases for $t>e$.
For $w=L-5c$ the same inequality follows from
$w'/w=(L-5)/(L-5c)$; its excess over one is decreasing here.

In particular the gap needed for the rectangle cutoff,
$L/2-10c-10\log w_3$, has derivative
$L/2-10-10w_3'/w_3>L/2-20.1>0$.
Its certified positive endpoint proves $D_0<D_R$ throughout.
The lower bound $\log D_R\ge L/2-10c$ also proves
$D_R>10^9$ throughout.  The small-conductor support gaps
$\kappa_1L-10c$ and $\kappa_1L-10\log w_q$ have positive
derivatives, so all removed exceptional primes remain below $z_*$.

The ratio of $\nu$ to $1/(2\cdot2.0452\cdot10c)$ is a constant
times $e^{-\delta c/2}/c$, hence decreases.
It is below one at $c=24.9$.  All small-exceptional thresholds are at
most $L^\delta$, so the common $\nu$ is admissible in every stage.
Moreover, for each of the three choices of $w$ above,
\begin{equation}\label{eq:nu-growth}
 \frac{d}{dc}\log(\nu w)
 =\frac{w'}w-\frac{\delta}{2}-\frac2c
 >1-\frac{29}{50}-\frac{20}{249}>\frac13.
\end{equation}
The certificate gives $\nu w>100$ at the smallest starting scale.
It follows that $\nu w>100$ on the entire half-line.

These inequalities also justify every maximum over prime-counting
endpoints in \eqref{eq:vformula}: for fixed $N$ and $\nu(N)$,
differentiate with respect to $x=\log y\ge w$.
The derivatives of $x^4e^{-\nu x}$,
$x^{14}e^{-x/2}$ and $x^{14}e^{-2x/3}$ have respective signs of
$4-\nu x$, $14-x/2$, and $14-2x/3$, all negative.
The first term $x^{-4}$ decreases too.  This proves the endpoint
majorant for all intermediate $y$, not just the endpoint $Y$.

For variation in $N$, the logarithmic derivative of
$Lw^4e^{-\nu w}/(1-\nu)$ is at most
$1+4(1.01)-\nu w/3<0$; the derivative of $(1-\nu)^{-1}$ only
helps.  The corresponding derivatives for
$Lw^{14}e^{-w/2}$ and $Lw^{14}e^{-2w/3}$ are negative by
\eqref{eq:wderivative}.  The derivative of $L/w^4$ is negative.
The multiplier in the second line of \eqref{eq:vformula} decreases,
and $L/\log(N^b/2)=1/(b-\log2/L)$ decreases.
Therefore $Lv(N^b/2)$, and hence $v(N^b/2)$, decrease.
In particular the certified inequality $v(X_q)<40/L$ holds forever.

\subsection{Rectangular and ordinary AP errors}

Every term of $m_N$ in \eqref{eq:mrect} decreases:
the first by the preceding paragraph;
the second because $t^{16}e^{-t}/\log t$ has negative derivative
for $t>16$; the third because
$d\log(L^5/w_3^{10})/dc=5-10w_3'/w_3<-5$;
the fourth because $5-\alpha L<0$ for
$\alpha=0.199,1/16$; and the last because $\log X_3$ increases.
After summing the ceiling bound, the right side of
\eqref{eq:box-error} is exactly a constant times
$m_N(A+B/L)$ for fixed positive constants $A,B$.
It therefore decreases.  This proves the global box assertion despite
the growing number of boxes.

For the AP coefficient, use the explicit formula in
\cref{sec:APformula}.  Let $w=L-5c$ and $h=w-5\log w$.
Besides \eqref{eq:wderivative}, $h'/h>w'/w>1$.
Each nonconstant summand of $L^2b(N)$ decreases.
For the power part of $\mathcal E$ its largest derivative is bounded by
$2+(13/2)w'/w-10h'/h< -3/2$.
The other power part has derivative $2-(7/2)w'/w< -3/2$.
The zero term $L^2w^2\log w\,e^{-\nu w}/(1-\nu)$ has logarithmic
derivative at most
$2+2(1.01)+1.01/\log w-\nu w/3<0$.
For every remaining exponential term, the highest polynomial degree
in $w$ is $15/2$ and its smallest exponential decay is $e^{-w/12}$;
$2+(15/2)(1.01)-w'/12<0$ bounds them.
The nonexponential term $L^2\log w/w^6$ has derivative
$2+(1/\log w-6)w'/w<0$.
Terms with $w^{-8}$ have still smaller derivatives, and the rational
multiplier defining $p$ in \cref{sec:APformula} decreases.
Finally the added $\omega(N)$ term in $L^2b(N)$ has logarithmic
derivative
\[
 \frac{L-5}{2}-L+6-10\frac{L-5}{L-5c}-\frac1c<0.
\]
Thus the endpoint check $L^2b(N)<2.21$ proves the claimed AP bound
for all $c\ge24.9$.  Inserting this into \eqref{eq:ordinary-AP-cost}
leaves a constant times
$(1/L+1/\log2)^2/L$, which decreases.

\subsection{Switched, local, finite and transfer errors}

The normalized high-conductor envelope in the structural certificate
is a positive polynomial prefactor of logarithmic derivative at most
$6$, times the sum of
\[
 \frac8{w_q^{10}},\qquad
 \sqrt{12}K_R e^{-L/12},\qquad
 \sqrt{12}K_R e^{-L/3},\qquad
 24e^{-0.0846L}.
\]
Here $K_R'/K_R<1$.  The first product has derivative less than
$6-10=-4$.  The other three have derivatives bounded respectively by
$7-L/12$, $7-L/3$, and $6-0.0846L$, all negative.
This accounts for every dyadic conductor interval and every $n$ block.
For the low-conductor term insert $v(X_q)\le40/L$ in
\eqref{eq:switch-low}.  After normalization its varying factor is
$(1+L)\log w_q/L^3$, with logarithmic derivative at most
$1+1.01/\log w_q-3<-1.9$.

The Perron norm bound follows from
$\operatorname{arsinh}(TL)<\log(2TL)+1
 =4L+c+\log2+1$.
At the endpoint this is below $2\pi L$, and the difference has
positive derivative since $2\pi L-(4L+1)>0$.
The normalized tail is bounded by
$e^{20}L^{20}e^{-1.5846L}$, whose derivative has sign
$20-1.5846L<0$.
The integer rough-sieve remainder $e^{20}L^{-9}$ decreases.
Hence the entire switched-error allocation remains valid.

The exceptional mass bound $B_E$ has logarithmic derivative
$2/(c\log(10c))-\delta(L-5)/(L-5c)<0$, and $L^{-2}$ decreases.
In \eqref{eq:transfer-cost}, $d_N$ and $h_0$ decrease term by term.
In \eqref{eq:local-cost}, $L^{-2}$ and $L^{-9}$ decrease, while
$LN^{-\kappa_1}$ has derivative of sign $1-\kappa_1L<0$.

For completeness, \eqref{eq:finite-cost} follows without a numerical
claim about a sum of astronomical integers.  Since $\log2>1/2$,
$W<2L$, $1+W<3L$, and $B<4N^{11/12}$ here.
After division by $4\Th>2N/L^2$, the five terms of
\eqref{eq:finite-envelope} are bounded by
\[
 64L^4N^{-1/12},\quad 12L^4N^{-1/3},\quad12L^2N^{-2/3},
 \quad486L^7N^{-11/12},\quad36450L^8N^{-1/12}.
\]
Their sum is at most $37024L^8N^{-1/12}
 <e^{30}L^{12}N^{-1/12}$.  Its logarithmic derivative is
$12-L/12<0$, and its endpoint logarithm is below $\log10^{-6}$.

This proves every analytic propagation used in \eqref{eq:assembled}.
The main integrals contain only fixed rational parameters, not varying
endpoints in $c$.  All remaining quantities are bounded by the
decreasing envelopes just proved.  Thus ceiling jumps, changes of
source primes, and changes of exceptional conductor are covered at
their two sides without assuming that the original discrete
expressions are continuous or monotone.

\appendix
\section{The AP coefficient used in the certificate}\label{sec:APformula}

This appendix specifies the BJS formulas used in \eqref{eq:AP}; it also
explains why no maximum or hidden choice of threshold is delegated to
an uncertified computation.  Fix $N$ and put $w=L-5c$, $\nu=\nu(N)$.
For $x\ge w$ define
\[
 \mathcal E_0(x)=
 \frac{4x^{9/2}}{(x-5\log x)^{10}}+\frac4{x^{11/2}}
 +\frac{18e^{-x/12}}{\sqrt x}+\frac52 e^{-x/6}x^{11/2}.
\]
The upper coefficient obtained from BJS's bound for the sum of
$|E_\psi(y;d,N)|$, with $x=\log y\in[w,L]$, is
\begin{align*}
 p_2={}&w^2\left[
 1.1(10\log w)\left(\frac{3.2\cdot10^{-8}}{w^8}
                       +\frac{e^{-\nu w}}{1-\nu}\right)
 +27\mathcal E_0(w)
 +\frac{e^{-w/2}}{2(\log2)w^8}+0.4w^3e^{-w}\right],\\
 p_1={}&p_2+\frac{0.67+2e^{-w/6}}{w^8},\\
 p={}&p_1\left(1+\frac1{L^2w^3}
                 +\frac1{(1-4/w)L}\right)+\frac{2.2}{L^2},\\
 b(N)={}&p+\frac{0.9e^{w/2-L}L^4}{w^{10}c}.
\end{align*}
These are BJS's $p_2,p_1,p,c_4$ bounds with a weaker common zero
bound.  They apply also to the nonexceptional and fixed full-conductor
relative sums, since their zero gap is larger.

To justify using the displayed endpoint for $p_2$, at fixed $N$ replace
$\log Q_1(x_1(N))$ by $10\log x$ and multiply the BJS bound by
$x^2/e^x$.  Every resulting term decreases for $x\ge w$:
for the zero term its logarithmic derivative is
$2/x+1/(x\log x)-\nu<0$;
for the leading $\mathcal E_0$ power term it is at most
$(13/2)/x-10/x<0$, since
$(1-5/x)/(x-5\log x)>1/x$.
The remaining powers decrease, and for the exponentials the largest
degree is $15/2$ while the smallest decay rate is $1/12$.
All threshold inequalities hold already at $w(N_0)$.
This proves the maximum over the finite interval $[x_1(N),N]$;
extending it to $x\ge w$ only enlarges the allowable majorant.

The passage $\psi\to\theta$ gives the displayed $p_1$; partial summation
from $x_1(N)$ to $N$ and Brun--Titchmarsh at the initial point give
the displayed $p$, including the $2.2/L^2$ term.
Finally removing primes dividing $N$ costs at most
$1.3841L/c$ per modulus.  There are fewer than
$0.65H_N$ square-free moduli by BJS's square-free counting bound.
Since $0.65\cdot1.3841<0.9$, their normalized cost is exactly bounded
by the last term of $b(N)$.
These steps use the proofs of BJS Lemmas 31, 34, 35, and 43, on their
stated square-free, coprime modulus ranges.  They do not extend the
relative result to an arbitrary additional tag.

\section{Interval integration and reproducibility}\label{sec:certificate}

The numerical environment is Python 3.12 with \cert{python-flint}
0.9.0, using Arb/ACB at 192-bit precision and one thread.
All input constants are integers, exact rational quotients or decimal
strings; no binary floating-point approximations enter the certificates.
The rational identities are checked with exact rational arithmetic.
Every adaptive integral must return a finite real enclosure with at
least 90 bits of relative accuracy; an undecided or false comparison
aborts execution.  The integration limit is $200000$ function
evaluations and depth $60$, with no fallback to point sampling.

Here are the analytic reductions used to avoid unverifiable nested
quadrature.  For $3\le s\le5$,
\[
 F(s)=\frac{2e^\gamma}{s}
 \left[1+\log(s-2)\log(s-1)+\operatorname{Li}_2(2-s)
                                      +\frac{\pi^2}{12}\right].
\]
Differentiate the bracket to obtain $\log(s-2)/(s-1)$ and use
its value $1$ at $s=3$; this proves the formula from the delay equation.
For $4\le s\le6$ and $5\le s\le7$, respectively,
\begin{align*}
 f(s)&=\frac{2e^\gamma\log3+\int_3^{s-1}F(u)\,du}{s},\\
 F(s)&=\frac1s\left[
 5F(5)+2e^\gamma\log3\,\log\frac{s-1}{4}
 +\int_3^{s-2}F(u)\log\frac{s-1}{u+1}\,du\right].
\end{align*}
The second identity follows by inserting the first in
$sF(s)=5F(5)+\int_5^sf(v-1)\,dv$ and reversing the finite integrals.
These identities supply every function value beyond the printed
elementary $s\le4$ range.

For $L_4$, reflect the symmetric triangle to half of its square and
then set $v=t+u$.  The one-dimensional weight is
\[
 \begin{cases}
 v^{-1}\log((v-\kappa_1)/\kappa_1),
       &2\kappa_1\le v\le\kappa_1+\kappa_2,\\
 v^{-1}\log(\kappa_2/(v-\kappa_2)),
       &\kappa_1+\kappa_2\le v\le2\kappa_2 .
 \end{cases}
\]
For $L_5$ the corresponding weight is
\[
 \begin{cases}
 v^{-1}\log((v-\kappa_2)(v-\kappa_1)/(\kappa_1\kappa_2)),
       &\kappa_1+\kappa_2\le v\le2\kappa_2,\\
 v^{-1}\log(\kappa_2(v-\kappa_1)/(\kappa_1(v-\kappa_2))),
       &2\kappa_2\le v\le\lambda .
 \end{cases}
\]
The formulas follow by integrating $1/(t(v-t))$ over the relevant
line section, whose antiderivative is
$v^{-1}\log(t/(v-t))$.
The $A_2,A_3$ integrals are split where their arguments equal $3$ and
$2$; the latter split is new because the retuned upper endpoints enter
the branch $1<s<2$.  The exceptional $L_5$ integral is split at exactly
$v=\theta_- -3\kappa_1$ and $v=\theta_- -2\kappa_1$,
not at the unshifted values.  The latter is the jump point of $\Flow$.
The other integrals are split wherever their delay-function formula
changes.  Thus no nonsmooth formula is passed across an unmarked
integration boundary.

The sole forward entry point is
\begin{center}
\cert{python c249_verify.py --check-tex}.
\end{center}
It executes both component certificates afresh, reconstructs the signed
rows of \eqref{eq:master}, checks every printed interval against the
new ball enclosure, sums the seven budgets, and writes
\cert{c249_final_ledger.json}.
That file records directions, Wu weights, normalization, complete balls,
engine precision and SHA-256 hashes of the scripts, this manuscript,
and the exact two external source files used.
The component JSON files are reproducible subordinate outputs, not
alternative ledgers.  Previous $23.36$ and diagnostic floating-point
ledgers are not inputs to this proof.

The numerical run is not a formal proof checker for the analytic
lemmas.  Their proofs and their hypotheses are the body of this
manuscript.  The role of the run is solely to enclose the finite
computations and certify the explicit starting inequalities that the
half-line argument propagates.

\begin{thebibliography}{9}
\bibitem{BJS}
M. Bordignon, D. R. Johnston and V. Starichkova,
\emph{An explicit version of Chen's theorem and the linear sieve},
arXiv:2207.09452, version 6.
\url{https://arxiv.org/html/2207.09452v6}.
\bibitem{Wu}
J. Wu, \emph{Chen's double sieve, Goldbach's conjecture and the twin
prime problem}, Acta Arithmetica \textbf{114} (2004), 215--273.
\url{https://arxiv.org/abs/0705.1652}.
\end{thebibliography}
\end{document}

